\documentclass[11pt]{article}
\usepackage[margin=1in]{geometry}
\usepackage{amsmath,amssymb,amsthm,mathtools,booktabs,array}

\usepackage{microtype}
\usepackage[colorlinks=true,linkcolor=blue,citecolor=blue,urlcolor=blue]{hyperref}
\usepackage{enumitem,fancyhdr}
\newtheorem{theorem}{Theorem}
\newtheorem{proposition}[theorem]{Proposition}
\newtheorem{corollary}[theorem]{Corollary}
\newtheorem{lemma}[theorem]{Lemma}
\theoremstyle{remark}\newtheorem{remark}[theorem]{Remark}
\newcommand{\R}{\mathbb R}
\newcommand{\T}{\mathsf T}
\newcommand{\diag}{\operatorname{diag}}
\newcommand{\Var}{\operatorname{Var}}
\newcommand{\eff}{\operatorname{Eff}}
\newcommand{\norm}[1]{\left\lVert#1\right\rVert}
\newcommand{\ip}[2]{\left\langle#1,#2\right\rangle}
\newcommand{\simplex}{\Delta_m^\circ}
\pagestyle{fancy}\fancyhf{}\fancyhead[L]{\small Delayed experimental design}
\fancyhead[R]{\small WORKING PAPER}\fancyfoot[C]{\thepage}
\setlength{\headheight}{14pt}
\setlength{\parskip}{3pt}
\setlist{nosep}
\newcommand{\internalnotice}{\begin{center}\small\textbf{RESEARCH WORKING PAPER --- 26 November 2026}\\Prepared with AI assistance and checked in separate AI reviews.\par Human author verification is pending.\end{center}}

\setlength{\emergencystretch}{1em}

\renewcommand{\thetheorem}{S\arabic{theorem}}
\renewcommand{\thesection}{S\arabic{section}}
\renewcommand{\theequation}{S\arabic{equation}}
\title{Supplement to\\Certifying Delayed Multiplicative Updates\\for Linear Experimental Design}
\author{}\date{26 November 2026}
\begin{document}
\maketitle\internalnotice
This supplement gives self-contained proofs, exact worked examples, an implementation contract, and a reproduction guide. All proved guarantees concern exact arithmetic and known finite linear regressors. Numerical checks illustrate the statements; they are not substitutes for the proofs.

\section{Model, value, and tangent gap}
Let $f_i=(a_i,z_i^\top)^\top\in\R^{p+1}$, $i=1,\ldots,m$, and assume the matrix $F$ of feature rows has full column rank. For a positive probability vector $w$, put
\[
 A_w=\sum_iw_iz_iz_i^\top,\quad b_w=\sum_iw_iz_ia_i,
 \quad c_w=\sum_iw_i a_i^2.
\]
Then $A_w\succ0$ and
\begin{equation}\label{eq:svalue}
 \theta_w=-A_w^{-1}b_w,\quad
 v(w)=c_w-b_w^\top A_w^{-1}b_w>0.
\end{equation}
Indeed $M_w=F^\top\diag(w)F\succ0$, and its Schur complement of $A_w$ is $v(w)$. The inverse block formula gives $(M_w^{-1})_{11}=1/v(w)$. For independent homoscedastic observations with exactly $Nw_i$ observations at site $i$, the information matrix is $NM_w/\sigma^2$ and the least-squares target variance is $\sigma^2/(Nv(w))$. This statement assumes integral counts and full rank; it makes no rounding claim.

Fix a positive cache $s$. Set $p_i=a_i+z_i^\top\theta_s$, $v_s=\sum_i s_i p_i^2$, and store a factorization of $A_s$. The cache normal equation is $\sum_i s_i z_i p_i=0$. At any positive current design $w$, let
\[
 q_w=\sum_iw_ip_i^2,\qquad r_w=\sum_iw_iz_ip_i.
\]
For $u=\theta-\theta_s$, the contrast energy is exactly
\begin{equation}\label{eq:sexpand}
 \sum_iw_i(a_i+z_i^\top\theta)^2
 =q_w+2r_w^\top u+u^\top A_wu.
\end{equation}
Completing the square proves
\begin{equation}\label{eq:sresid}
 \theta_w=\theta_s-A_w^{-1}r_w,\quad
 v(w)=q_w-r_w^\top A_w^{-1}r_w.
\end{equation}
In particular
\[
 D(w\mid s):=v(s)+\nabla v(s)^\top(w-s)-v(w)
 =r_w^\top A_w^{-1}r_w
 =\sum_iw_i(p_i-p_{w,i})^2.
\]
To justify the gradient without an unproved differentiation step, note that the normal equations give $\sum_i w_i z_i p_{w,i}=0$. Differentiating~\eqref{eq:svalue} or~\eqref{eq:sexpand} therefore cancels the terms involving $d\theta_w$, leaving $dv=\sum_i p_{w,i}^2dw_i$. Smoothness follows from $A_w\succ0$. The ambient gradient is defined for the homogeneous extension to positive weights, so it has no arbitrary additive constant from simplex coordinates.

\section{Full proof of the safeguard}
Suppress the subscript $w$ on $q,r$. Define
\[
 T_i=w_i p_i^2/q,\quad d_i=T_i-w_i,\quad
 h=\sum_i d_i z_i p_i,\quad B=\sum_i d_i z_i z_i^\top.
\]
Since $q>0$, $T$ is a probability vector and $\sum_i d_i=0$.
Let
\begin{align*}
 Q&=r^\top A_s^{-1}r,&H&=h^\top A_s^{-1}h,&
 \alpha&=\min_i w_i/s_i,\\
 \rho_i&=d_i/w_i=p_i^2/q-1,&
 \gamma&=\max_i|\rho_i|,&\zeta&=\max_i(-\rho_i),\\
 g&=\sum_i d_i p_i^2
 =\frac{\sum_iw_ip_i^4-q^2}{q},&&&
 G&=g-\frac{2\sqrt{HQ}+\zeta Q}{\alpha},\\
 K&=\frac{(\sqrt H+\gamma\sqrt Q)^2}{\alpha}.&&&&
\end{align*}
The zero weighted mean of $\rho$ and $\rho_i\ge-1$ imply $0\le\zeta\le1$. Both $Q,H$ are nonnegative and $g\ge0$.

\begin{theorem}\label{thm:smain}
For $0\le t<1$, all updated weights $w_i+td_i$ are positive, and
\[
 v(w+td)-v(w)\ge tG-\frac{t^2K}{1-t\zeta}.
\]
If $G>0$, choose $t=\min\{1/2,G/(4K)\}$ for $K>0$ and $t=1/2$ for $K=0$. Then $v(w+td)-v(w)\ge tG/2>0$.
\end{theorem}
\begin{proof}
Let $A=A_w$ and $e=A^{-1}r$. These current objects are analysis variables, not algorithm inputs. The componentwise bounds on $w_i/s_i$ imply $A\succeq\alpha A_s$ and hence
\begin{equation}\label{eq:snorms}
 r^\top A^{-1}r\le Q/\alpha,
 \qquad h^\top A^{-1}h\le H/\alpha.
\end{equation}
The componentwise bounds on $d_i/w_i$ imply
\begin{equation}\label{eq:sB}
 -\zeta A\preceq B\preceq\gamma A,\qquad
 \norm{A^{-1/2}BA^{-1/2}}_2\le\gamma.
\end{equation}
For the second inequality, the negative eigenvalues have magnitude at most $\zeta\le\gamma$.
Write $P_i=p_i-z_i^\top e$ for the unknown current optimal contrast.
The true directional derivative equals
\[
 L:=\sum_i d_iP_i^2=g-2h^\top e+e^\top Be.
\]
Cauchy--Schwarz in the $A^{-1}$ metric and~\eqref{eq:snorms} give
$|h^\top e|\le\sqrt{HQ}/\alpha$.
Also~\eqref{eq:sB} gives $e^\top Be\ge-\zeta Q/\alpha$.
Therefore $L\ge G$.

Now $A_t=A+tB\succeq(1-t\zeta)A\succ0$.
The residual of $P$ under the new design is
\[
 \sum_i(w_i+td_i)z_iP_i=t(h-Be)=:tb.
\]
Completing the square at the coefficient $\theta_w$ gives the exact, finite-step identity
\begin{equation}\label{eq:sexactchange}
 v(w+td)-v(w)=tL-t^2b^\top A_t^{-1}b.
\end{equation}
There is no omitted Taylor remainder. By~\eqref{eq:sB},
\[
 \norm b_{A^{-1}}
 \le\norm h_{A^{-1}}+\norm{Be}_{A^{-1}}
 \le\sqrt{H/\alpha}+\gamma\sqrt{Q/\alpha}=\sqrt K.
\]
Thus $b^\top A_t^{-1}b\le K/(1-t\zeta)$, proving the first bound.
For the recommended step, $t\le1/2$ and so $1-t\zeta\ge1/2$.
If $K>0$, $tK\le G/4$; if $K=0$, the whole remainder vanishes.
In either case the claimed gain follows. Positivity follows from
$w_i+td_i=w_i(1+t\rho_i)\ge(1-t)w_i>0$.
\end{proof}

\paragraph{Boundary and zero cases.}
Some $p_i$ may be zero; no division by $p_i$ occurs and damping retains that site's positive weight. If $Q=0$, then $r=0$ and the cached coefficient is also current-optimal. If $K=0$, the proof forces $b=0$ and the objective is affine on the segment. If $d=0$, then $g=G=0$, so no false strict gain is claimed. If $g=0$, positivity of all $w_i$ forces every $p_i^2$ to be equal and hence $d=0$. Stale zero motion does not itself establish that the current design is optimal; the optimality conclusion requires a current cache. Singular caches, nonestimable targets, and floating-point error bounds are not covered.

\begin{corollary}[Local nonvacuity with a positive lower step]\label{cor:slocal}
Fix a positive cache $s$ with nonconstant squared sensitivities. Put $g_s=g(s)>0$ and $K_s=K(s)$. There is a relative neighborhood of $s$ on which
\[
 G(w)\ge g_s/2,\qquad K(w)\le K_s+g_s,
\]
and the selected step satisfies
\[
 t(w)\ge t_0:=\min\{1/2,g_s/[8(K_s+g_s)]\}>0,
 \qquad v(w+t(w)d(w))-v(w)\ge t_0g_s/4.
\]
\end{corollary}
\begin{proof}
At $w=s$ the cache normal equation gives $Q=0$, $\alpha=1$ and
$G=g_s=\Var_s(p_i^2)/v_s>0$.
The finite sums, ratios with positive denominators, square roots and finite maxima defining $G,K$ are continuous. Choose a neighborhood satisfying the displayed inequalities and apply Theorem~\ref{thm:smain}. If $K=0$, the selected step is $1/2$, which satisfies the same bound.
\end{proof}
The neighborhood is not asserted to contain the next point after an accepted step. In particular this corollary cannot be iterated to obtain an unproved number of cache reuses.

\section{Efficiency and value intervals}
For a positive candidate $u$, define $q_u,r_u,Q_u$ as above, and
$\alpha_u=\min_i u_i/s_i$, $\beta_u=\max_i u_i/s_i$.
Positive-definite comparison gives
\begin{equation}\label{eq:sinterval}
 q_u-Q_u/\alpha_u\le v(u)\le q_u-Q_u/\beta_u.
\end{equation}
For the positive dual lower bound, put $\ell_i=s_i p_i/v_s$.
The normal equation implies $\sum_i\ell_i z_i=0$, and
\[
 \sum_i\ell_i a_i
 =v_s^{-1}\sum_i s_i p_i(p_i-z_i^\top\theta_s)=1.
\]
Consequently $\sum_i\ell_i(a_i+z_i^\top\theta)=1$ for every $\theta$.
Weighted Cauchy--Schwarz gives
\[
 v(u)\ge L_{\rm dual}(u):=
 \left(\sum_i\ell_i^2/u_i\right)^{-1}>0.
\]
Define $L_s(u)=\max\{q_u-Q_u/\alpha_u,L_{\rm dual}(u)\}$ and
$U_s(u)=q_u-Q_u/\beta_u$.
Then $L_s(u)\le v(u)\le U_s(u)$, so $L_s(u)\ge U_s(w)$ certifies a nondecrease.
The maximum cached sensitivity $V_s=\max_i p_i^2$ bounds $v_*$ from above because the cached contrast is feasible for every design. Hence
\begin{equation}\label{eq:seff}
 \eff(u)\ge L_s(u)/V_s.
\end{equation}
Combining this with an accepted damped step yields the possibly sharper
$\eff(w+td)\ge[L_s(w)+tG/2]/V_s$.
All these bounds are deterministic design-value bounds, not confidence intervals for random observations.

\paragraph{Incremental evaluation.}
After computing $A_s^{-1}r$ and $A_s^{-1}h$, form
$C=r^\top A_s^{-1}h$.
Along the chosen line,
\[
 q_{w+td}=q+tg,\quad r_{w+td}=r+th,\quad
 Q_{w+td}=Q+2tC+t^2H.
\]
The new weight ratios and reciprocal dual sum take $O(m)$ work. Thus evaluating the interval after the certified step requires no third cached solve. A completely unrelated candidate generally requires an additional residual contraction and cached solve.

\section{Exact diagnostic examples}
\subsection{A reachable efficiency collapse}
For $x=(1,b)$, $b>1$, and first weight $u$, direct minimization gives
\[
 \theta_u=-\frac{u+b(1-u)}{u+b^2(1-u)},\qquad
 v(u)=\frac{(b-1)^2u(1-u)}{u+b^2(1-u)}.
\]
The maximizing first weight is $b/(b+1)$, with
$v_*=(b-1)^2/(b+1)^2$.
At $s=(1/2,1/2)$,
\[
 p_1=\frac{b(b-1)}{1+b^2},\quad
 p_2=-\frac{b-1}{1+b^2}.
\]
Thus the fresh first proposal is $w=(b^2,1)/(1+b^2)$ and the next stale full proposal is $u=(b^4,1)/(1+b^4)$.
Substitution gives
\[
 v(s)=v(w)=\frac{(b-1)^2}{2(1+b^2)},\quad
 v(u)=\frac{b^2(b-1)^2}{(1+b^4)(1+b^2)},
\]
proving the efficiency-collapse claim in the main text.
At $b=2$, $s=(1/2,1/2)$, $w=(4/5,1/5)$,
$u=(16/17,1/17)$, $v(w)=1/10$, and $v(u)=4/85$.
The directional certificate rejects this harmful proposal: its exact $G$ is $-828/10625$.

\subsection{Arbitrarily small drift with the wrong sign}
Take two feature rows $(1,1)$ and $(1,-1)$ and let their first weights be $y$ at the cache and $x$ currently. Then
$v(x)=4x(1-x)$ and $T_y(x)=x(1-y)^2/[x(1-y)^2+(1-x)y^2]$.
Choose $y=1/2-\epsilon$, $x=1/2+\epsilon$, where $0<\epsilon<1/2$.
This pair is reachable by a fresh full update, since $T_y(y)=1-y=x$.
The following exact quantities make the obstruction quantitative:
\begin{align*}
 \norm{w-s}_1&=4\epsilon,&
 \alpha&=(1-2\epsilon)/(1+2\epsilon),&
 D(w\mid s)&=16\epsilon^2,\\
 q&=1+12\epsilon^2,&
 d_1&=\frac{2\epsilon(1-4\epsilon^2)}{1+12\epsilon^2},&
 g&=\frac{16\epsilon^2(1-4\epsilon^2)}{1+12\epsilon^2}.
\end{align*}
The true slope is $\nabla v(w)^\top d=-g<0$.
Every $0<t\le1$ moves farther to the right of the maximizer and decreases $v$.
Thus an absolute tolerance on mismatch or relative weight drift is insufficient without accounting for an ascent margin. This does not contradict a margin-dependent local neighborhood.

\subsection{A certified second step without a new factorization}
Take $x=(1,2,3)$ and $s=(1/3,1/3,1/3)$.
Then $A_s=14/3$ and $p=(4,1,-2)/7$.
The first safe step from the cache has $t_0=1/2$ and produces
$w=(23,8,11)/42$.
At this genuinely stale design, direct rational arithmetic gives
\[
 q=10/49,\quad g=156/1715,\quad r=1/7,\quad h=67/245,
 \quad\alpha=4/7,\quad\gamma=\zeta=9/10.
\]
Since there is one nuisance parameter, the square-root products simplify exactly:
$\sqrt{HQ}=|hr|/A_s$ and
$K=(|h|+\gamma|r|)^2/(A_s\alpha)$.
Consequently
\[
 G=1503/27440>0,\qquad K=116427/1920800,\qquad
 t=17535/77618<1/2.
\]
Theorem~\ref{thm:smain} certifies this second update using the original cache. The accompanying exact-arithmetic program independently recomputes the new weights and checks the actual rational gain against $tG/2$.

\subsection{A full proposal certified by interval separation}
Take $x=(1,2,3)$ and $s=(1,7,1)/9$.
Here $p=(10,1,-8)/19$, $v_s=1/19$,
$w=(100,7,64)/171$, and $u=(10000,7,4096)/14103$.
Using~\eqref{eq:sinterval} and the dual bound gives
\[
 U_s(w)=1481081/6516050,
 \qquad L_s(u)=6931/29773.
\]
Their difference is positive by integer cross multiplication.
The exact target values are $v(w)=743/3344$ and
$v(u)=4553852/18369941$.
For this pair the directional lower bound has $G<0$, so that test rejects while the full interval test accepts. The example demonstrates conservative rejection, not a contradiction between certificates.


\section{Two limits of reliable local updates}
\subsection{Vanishing relative mismatch and efficiency collapse}
At nodes $(1,K)$ with $K>1$, parameterize the positive design by first-weight odds $u$. The normal equation and value are
\[
 \theta_u=-\frac{u+K}{u+K^2},\quad
 P_u(1)=\frac{K(K-1)}{u+K^2},\quad
 P_u(K)=-\frac{u(K-1)}{u+K^2},\quad
 v(u)=\frac{u(K-1)^2}{(1+u)(u+K^2)}.
\]
These formulas show both that the nodal values never vanish and that an old contrast from odds $a$ updates current odds $b$ to $bK^2/a^2$.
The optimum is at $u=K$, as the exact identity
\[
 (1+u)(u+K^2)-u(K+1)^2=(u-K)^2
\]
proves $v(u)\le(K-1)^2/(K+1)^2$ with equality there.
For any positive $a,b$, the nuisance information is $A_b=(b+K^2)/(1+b)$, so the square-completion gap gives
\[
 \frac{D(\pi(b)\mid\pi(a))}{v(b)}
 =\frac{K^2(a-b)^2}{b(a+K^2)^2}.
\]
Now take an integer $n\ge2$, put $K=n^2$, and use identical initial odds $u_0=u_1=n^3$. The successive odds are
\[
 n^3,n^3,n,n^{-1},n,n^7.
\]
At the step from $u_4=n$ using $u_3=n^{-1}$, substitution gives
\[
 d_n=\frac{n^3(n^2-1)^2}{(n^5+1)^2},\qquad
 \frac{v(u_5)}{v(u_4)}=\frac{n^3(n+1)}{n^7+1}.
\]
Both quantities tend to zero as $n\to\infty$. The current efficiency equals
\[
 E_n=\frac{(n^2+1)^2}{(n+1)^2(n^2-n+1)},\qquad
 1-E_n=\frac{n(n-1)^2}{(n+1)^2(n^2-n+1)}\to0.
\]
Hence the asserted geometry-independent lower bound cannot hold. This proof controls exact compatible history and exact statistical efficiency; it does not turn this elementary two-node mechanism into a general lower bound for all delayed algorithms.

\subsection{Proof that refresh only upon failure can stall}
Let the two feature rows be $(1,-1)$ and $(1,1)$, cache $s=(1/4,3/4)$, and initial current design $w=s$. Write $x$ for the current first weight. The cache has $A_s=1$, $\theta_s=-1/2$, $p=(3/2,1/2)$, and the value is $v(x)=4x(1-x)$ with optimum $1$.
For $1/4\le x<1/2$, put
\[
 d=\frac{8x(1-x)}{1+8x},\quad R=2x-1/2,\quad
 \alpha=\frac{4(1-x)}3,\quad\gamma=d/x,\quad\zeta=d/(1-x).
\]
The direction is $(d,-d)$, $r=-R$, $h=-2d$, $Q=R^2$, $H=4d^2$, and $g=2d$. Direct substitution into the certificate gives
\[
 G=dZ(x),\quad Z(x)=\frac{80x^2-160x+53}{16(1-x)^2},\qquad
 K=\frac{d^2(4-1/(2x))^2}{\alpha}.
\]
Let $\rho=1-3\sqrt{15}/20$ and $\rho_+=1+3\sqrt{15}/20$.
Then $G>0$ for $1/4\le x<\rho$. The selected positive increment obeys
\begin{align*}
 td&\le\frac{\alpha Z(x)}{4(4-1/(2x))^2}\\
   &=(\rho-x)\frac{5(\rho_+-x)}{3(1-x)(4-1/(2x))^2}.
\end{align*}
On this interval the last factor is at most
\[
 C=\frac{5(5+\sqrt{15})}{12\sqrt{15}}<1.
\]
Indeed $\rho_+-x\le\rho_+-1/4$, $1-x\ge1-\rho$, and $4-1/(2x)\ge2$. The strict inequality $C<1$ is equivalent to $25<7\sqrt{15}$, whose square is $625<735$.
Thus every iterate strictly increases yet remains below $\rho$, and $G$ stays positive at every finite step: the cache never refreshes.
If the increasing sequence had limit $\ell<\rho$, continuity would give positive limiting $G,K,d$ and hence a positive limiting increment, a contradiction. Its limit is exactly $\rho$.
Finally
\[
 v(\rho)=4\rho(1-\rho)=3\sqrt{15}/5-27/20<0.99.
\]
The last strict inequality also follows by squaring $\sqrt{15}<39/10$.
This proves failure of the stated refresh-on-failure rule, not of every possible refresh policy.
Fresh half-power weighting has initial unnormalized masses
$(1/4)(3/2)=(3/4)(1/2)=3/8$ and therefore reaches the optimal half-half design immediately.

\subsection{A nonempty zero-curvature branch}
At positive nodes $(1,3,8)$, let $s=w=(15,14,6)/35$.
The cache is $A_s=15$, $\theta_s=-1/5$, and $p=(4,2,-3)/5$.
Then $r=h=0$, $Q=H=K=0$, while $g=G=3/25$.
The defined branch $t=1/2$ yields weights $(39/70,7/25,57/350)$, and the value increases from $2/5$ to $23/50$.
Thus the exact gain is $3/50=tG$. Omitting the zero-$K$ branch would leave a real admissible case undefined.
The supplied independent rational check also verifies the small-drift example, using 36 scalar offline solves; none are algorithm inputs.

\section{Algorithm and actual accounting}
The following protocol implements the proved damped rule.
\begin{enumerate}[label=\arabic*.]
\item At a positive cache $s$, form and factor $A_s$, solve for $\theta_s$, and cache $p_i,v_s$.
\item At the current positive $w$, form $q,d,r,h$ and all weight ratios. Use the cached factor for $A_s^{-1}r$ and $A_s^{-1}h$.
\item Evaluate $G,K$. If $G>0$, use the stated step and retain the cache. Optionally compute the incremental efficiency certificate.
\item If $G\le0$, count the rejected certificate. A refresh recomputes the cache at $w$ and pays the entire matrix-formation and factorization cost. If the refreshed sensitivities are constant, the current design is optimal. Otherwise the refreshed gate admits a positive step.
\end{enumerate}
A separately configured full-proposal interval test is valid if the proposal is positive. It must be evaluated and charged explicitly; it is not assumed in the damped protocol. The preceding stall proof shows that this exact refresh rule can fail to reach a requested efficiency even while every finite step is accepted. The relative-margin rule with a refresh-count cap is analyzed separately below; it is not the preceding refresh-on-failure rule.

For dense features the refresh costs $O(mp^2+p^3)$, while a stale gate costs $O(mp+p^2)$, including two solves against a cached factor. A triangular solve uses both forward and backward substitution; it is not counted as a scalar operation. Rejected gates, fresh retries, and true-value audits all have nonzero cost. The 240-case batch additionally computes 240 matrix condition estimates, generally involving singular-value calculations; these are separately charged offline diagnostics. Known feature storage is $O(mp)$; cache storage is $O(p^2+m)$.
Low-dimensional examples do not test the potential asymptotic saving. Fair performance comparisons would also include current-gradient damping, one or two preconditioned corrections using the same cached factor, maintained inverse or low-rank updates, and an LP-based finite $c$-design solver. This round does not claim to have completed that benchmark.


\section{A specialization at the same cached-solve budget of the classical surrogate argument}
The following comparison is included as prior-principle attribution, not as a second innovation. Write $P_i=p_i-z_i^\top e$ for the exact current contrast and $V=v(w)>0$. Its unbiased estimator coefficients $b_i=w_iP_i/V$ satisfy $\sum_i b_i f_i=e_1$. The classical covariance bound gives, for any positive design $u$,
\[
 v(u)\ge\left(\sum_i\frac{b_i^2}{u_i}\right)^{-1}.
\]
At $u=w$ this bound has equality. At $u=w+td$, with $\rho_i=d_i/w_i$, let
\[
 C_t=\sum_i\frac{w_i\rho_iP_i^2}{1+t\rho_i}.
\]
The surrogate denominator is exactly
\[
 \sum_i\frac{b_i^2}{w_i+td_i}
 =\frac{V-tC_t}{V^2}>0.
\]
Thus if $C_t\ge0$,
\begin{equation}\label{eq:syugain}
 v(w+td)-V\ge\frac{VtC_t}{V-tC_t}\ge tC_t.
\end{equation}
The exact current $P$ is not computed by the following bound. Since
\[
 C_t=L-t\sum_i\frac{w_i\rho_i^2P_i^2}{1+t\rho_i},\qquad L\ge G,
\]
put $R_2=\sum_iw_i\rho_i^2p_i^2$ and use the weighted triangle inequality:
\[
 \left(\sum_iw_i\rho_i^2P_i^2\right)^{1/2}
 \le\sqrt{R_2}+\gamma\left(\sum_iw_i(z_i^\top e)^2\right)^{1/2}
 \le\sqrt{R_2}+\gamma\sqrt{Q/\alpha}.
\]
Therefore, with $B_{\rm sur}=(\sqrt{R_2}+\gamma\sqrt{Q/\alpha})^2$,
\[
 C_t\ge G-\frac{tB_{\rm sur}}{1-t\zeta}.
\]
If $G>0$, choosing $t=\min\{1/2,G/(4B_{\rm sur})\}$ when $B_{\rm sur}>0$, and $t=1/2$ otherwise, ensures $C_t\ge G/2$. Equation~\eqref{eq:syugain} yields the same certified form of improvement $tG/2$.

The observable $G$ still uses the two old-factor quadratic forms $Q,H$; $R_2$ needs only a scalar pass over already available weights and contrast values. No new current solve or feature contraction is needed. This proves a safe-existence reduction at the same cached-solve budget to the classical auxiliary-function principle, augmented by the explicit residual slope bound. It does not assert that this alternative step equals the step based on $K$, that the original step always satisfies the surrogate decrease condition, or that all detailed formulas appeared in the cited source. It does show why a correct one-step formula and cache reuse alone do not establish a new algorithmic principle.

\section{Boundary-safe convergence of the capped margin protocol}
This section proves the new global statement without a minimum-weight assumption. The original uncapped relative-margin protocol is separate and fails in general by the later explicit boundary construction. All reasoning is in exact arithmetic.

\subsection{The classical finite-interpolant bound}
Let $Z$ have rows $z_i^\top$, and let $\mathcal I$ be the finite collection of $p$-row sets with nonsingular $Z_I$. Full rank ensures $\mathcal I$ is nonempty. Put
\[
 \theta_I=-Z_I^{-1}a_I,\qquad
 \omega_I(w)=\frac{\det(Z_I)^2\prod_{i\in I}w_i}{\det(Z^\top\diag(w)Z)}.
\]
Cauchy--Binet and Cramer's rule give the classical identity
\begin{equation}\label{eq:volume}
 \sum_{I\in\mathcal I}\omega_I(w)=1,\qquad
 \theta_w=\sum_{I\in\mathcal I}\omega_I(w)\theta_I.
\end{equation}
For completeness, replace column $j$ of $Z$ by $-a$ to obtain $Z^{(j)}$. Cramer's numerator for $(\theta_w)_j$ is $\det(Z^\top\diag(w)Z^{(j)})$. Its Cauchy--Binet expansion is
$\sum_I\det(Z_I)\det(Z_I^{(j)})\prod_{i\in I}w_i$.
Singular terms vanish, and each other term equals
$\det(Z_I)^2(\theta_I)_j\prod_{i\in I}w_i$.
Dividing proves~\eqref{eq:volume}. Thus no general-position assumption is needed.

This is the weighted least-squares convex-hull property of Ben-Tal and Teboulle~\cite{bental1990}, also the square-size specialization of the volume-sampling expectation identity in Derezi\'nski and Warmuth~\cite{derezinski2018}, Theorem~5 and Appendix~B. The identity is background, not a new result. It implies the finite constant
\begin{equation}\label{eq:uniformL}
 L=\max_{I\in\mathcal I,\,1\le i\le m}(a_i+z_i^\top\theta_I)^2>0,
 \qquad (a_i+z_i^\top\theta_w)^2\le L
\end{equation}
for every positive $w$. For $p=0$ the unique empty interpolant gives the same assertion directly. Computing $L$ is not part of the algorithm: it is a proof constant and may involve many subsets. The bound does not control $\|A_w^{-1}\|$ or impose a floor on the weights.

The closed-simplex optimum exists. Indeed $v$ is an upper-semicontinuous infimum of continuous affine functions of $w$, and is bounded on the compact simplex by the objective at $\theta=0$. The strictly positive starting value $v_0$ is positive by full rank. The gain guarantee ensures $v_0\le v(w_k)\le v_*$. At every cache $p_i^2\le L$, whence $v_*\le L$ and $q_k\ge v(w_k)\ge v_0$.

Giannelli et al.~\cite{giannelli2024}, Theorem~1 (published page 54), give the same finite-interpolant identity: set their collocation matrix to $Z$ and data vector to $-a$. Positive weights and full column rank hold here. This additional direct reference neither changes the classical attribution nor implies convergence of fits at a singular limit.

\subsection{Precise algorithm and the smaller-step guarantee}
The initial exact cache has counter $j=0$. If the current cache fails $g>0$ and $G\ge g/2$, replace it by the exact current fit and increment $j$. Stop only if this refreshed fit has $g=0$. Otherwise take
\[
 t=\min\{1/2,1/(j+1),G/(4K)\},
\]
with the last quantity omitted for $K=0$, and set $w^+=w+td$. A nonconstant fresh fit has $Q=0,G=g>0$, so a refreshed nonoptimal state always has a step. In practice initialization can check $g=0$ before a redundant refresh; this changes no accepted trajectory or proof.
For every chosen positive $t$, $w_i^+\ge w_i/2$. Since $t\le1/2$ and $tK\le G/4$, the original remainder bound gives
\begin{equation}\label{eq:smallerstep}
 v(w^+)-v(w)\ge tG-2t^2K\ge tG/2\ge tg/4.
\end{equation}
The zero-$K$ case is immediate. Finite termination occurs at an optimal fresh fit because equal positive-design sensitivities bound all design values above by the current value. An old cache with $g=0$ is not alone a reason to certify an arbitrary current design: the refresh instruction is essential.

\subsection{Finitely many refreshes}
Suppose there is an infinite sequence of accepted steps but only finitely many refreshes. Fix its last cache. Set $\lambda_i=p_i^2$, $M=\max_i\lambda_i$ and $m_0=\min_i\lambda_i$. If all $\lambda_i$ are equal, the protocol stops at that fresh cache; hence this infinite branch has $M>m_0\ge0$ and $M>0$.
The cap $c=\min\{1/2,1/(j+1)\}$ is now a fixed positive number.

If $\sum_k t_k<\infty$, the inequality $\log(1-t)\ge-2t$ for $t\le1/2$ implies
$w_{k,i}\ge w_{b,i}\exp(-2\sum_{k\ge b}t_k)>0$ uniformly after the last refresh $b$. Thus $\alpha$ is bounded below. The finite fixed cached data and $q\ge v_0$ bound $\gamma,Q,H,K$ above. Distinct $\lambda_i,\lambda_l$ give
\[
 g=\frac{\sum_{i<l}w_iw_l(\lambda_i-\lambda_l)^2}{q}\ge\underline g>0.
\]
The margin yields $G\ge\underline g/2$. Together with the fixed cap and bounded $K$, this gives a uniform positive lower bound on $t$, contradicting its summability. This conditional lower weight bound is derived from the contradiction hypothesis; it is not imposed on the actual run. Consequently $\sum_k t_k=\infty$.

Choose a maximum-score index $l$. For any $i$ with $\lambda_i<M$,
\[
 \frac{w_i^+/w_l^+}{w_i/w_l}
 =1-\frac{t(M-\lambda_i)}{q+t(M-q)}
 \le\exp\{-t(M-\lambda_i)/M\}.
\]
The denominator lies in $(0,M]$. Hence all lower-score mass vanishes, while ratios among maximum-score weights stay constant. In particular $q\to M$, $g\to0$, and
$\zeta=1-m_0/q\to1-m_0/M>0$.
The margin gives
\[
 0\le D(w\mid s)\le Q/\alpha\le g/(2\zeta)\to0.
\]
It follows that $v(w)\to M$. The fixed contrast bounds every design value by $M$, while every iterate has value at most $v_*$, so $M=v_*$. Indeed the weights converge to the normalized restriction of $w_b$ to the maximum-score indices. Since $\alpha\le1$, also $Q\to0$ and $r\to0$. At the limiting, possibly singular, design $w_\infty$,
\[
 \sum_i w_{\infty,i}(p_i+z_i^\top u)^2
 =M+u^\top A_{w_\infty}u\ge M,
\]
showing $v(w_\infty)=v_*$ without invoking boundary inverse continuity.

\subsection{Infinitely many refreshes}
At a fresh cache write $\rho_i=p_i^2/q-1$. The weighted projection contraction gives
\begin{align*}
 H&=\left\|A_w^{-1/2}\sum_iw_i\rho_i z_i p_i\right\|^2
 \le\sum_iw_i\rho_i^2p_i^2
 \le L\sum_iw_i\rho_i^2=L g/q.
\end{align*}
Here $Q=0$, $G=g$ and $K=H$. Thus the original fresh safe step is at least $\tau=v_0/(4L)>0$, including $H=0$. The capped fresh step at counter $j$ is therefore at least $\min\{\tau,1/(j+1)\}$. If refreshes are infinite, these steps alone have divergent sum. Moreover the current counter tends to infinity, so all accepted steps satisfy $t_k\to0$. Neither fact assumes a bound on epoch length.

For any epoch beginning at a fresh state $b$, hold its residual vector fixed in defining $q$. Then $q_b=v(w_b)$ and
$q_{k+1}-q_k=t_kg_k$ throughout that epoch. Summing~\eqref{eq:smallerstep} shows
\begin{equation}\label{eq:epoch-supp}
 0\le D_k=q_k-v(w_k)
 \le3\{v(w_k)-v(w_b)\}\le3\{\bar v-v(w_b)\},
\end{equation}
where $\bar v=\lim_k v(w_k)$ exists by monotonicity and boundedness. At a refresh $D=0$; at every other sufficiently late state its epoch start also tends to infinity. Therefore $D_k\to0$ and $q_k\to\bar v$. This argument does not require $g$ at the increasingly damped refreshes to tend to zero.

Fix any optimal $u_*\in\Delta_m$, used only in the proof. For each current cached contrast,
$\sum_i u_{*,i}p_{k,i}^2\ge v_*$ by the variational definition. Let
\[
 \rho_{k,i}=p_{k,i}^2/q_k-1,\quad
 R=\max\{1,L/v_0\},\quad
 \mathcal D_k=\sum_{i:u_{*,i}>0}u_{*,i}\log\frac{u_{*,i}}{w_{k,i}}.
\]
Then $-1\le\rho_{k,i}\le R$, $|\rho_{k,i}|\le R$, and the relative entropy $\mathcal D_k$ is finite and nonnegative. Suppose $\bar v<v_*$. Since $q_k\to\bar v$ and $q_k\le L$, eventually
\[
 \sum_i u_{*,i}\rho_{k,i}\ge\frac{v_*-q_k}{q_k}
 \ge\delta:=\frac{v_*-\bar v}{2L}>0.
\]
For $x\ge-1/2$, $\log(1+x)\ge x-x^2$: the derivative of their difference is $x(1+2x)/(1+x)$, negative on $[-1/2,0]$ and positive on $[0,\infty)$, with difference zero at zero. The weight update and $t_k\rho_{k,i}\ge-1/2$ now yield
\begin{align*}
 \mathcal D_{k+1}-\mathcal D_k
 &=-\sum_i u_{*,i}\log(1+t_k\rho_{k,i})\\
 &\le-t_k\sum_i u_{*,i}\rho_{k,i}
       +t_k^2\sum_i u_{*,i}\rho_{k,i}^2\\
 &\le-\delta t_k+R^2t_k^2.
\end{align*}
Eventually $t_k\le\delta/(2R^2)$, so each difference is at most $-\delta t_k/2$. Summing contradicts $\mathcal D_k\ge0$ because $\sum_k t_k=\infty$. Hence $\bar v=v_*$, completing the global proof. This is a classical diminishing-step entropy mechanism, within the established multiplicative-weights framework~\cite{freund1997}, Section~2 and Theorem~2 with its admissible linear multiplier; the design-specific obligations are the epoch error control, the uniform residual bound, and divergence of the refresh-step sum.

The proof also covers the capped surrogate step, replacing $K$ by $B_{\rm sur}$. The previous surrogate theorem supplies the identical smaller-step gain. At a fresh cache, $B_{\rm sur}=R_2\le Lg/q$. In the finite-refresh summability contradiction, an interior weight bound makes $B_{\rm sur}$ bounded as well. The remaining arguments are unchanged. Therefore this qualitative result establishes no separation from the classical surrogate route.

\subsection{What the uncapped argument does and does not imply}
The uncapped rule $t=\min\{1/2,G/(4K)\}$ with the same relative margin shares the finite-refresh result. At infinitely many fresh steps it has a uniform positive step lower bound, so their $g$ tends to zero by summable objective gains. The epoch argument still gives $D\to0$ and $q\to\bar v$. These facts alone do not imply global optimality, because $t_k\to0$ need not hold. The later actual-protocol counterexample shows that universal convergence is false.
An exact illustration of the weakness of weighted stationarity uses fixed rows
$(1,0),(2,1),(2,-1)$ and weights
$(1-\epsilon,\epsilon/100,99\epsilon/100)$.
The exact fit is $\theta=49/25$, $A=\epsilon$, and
\[
 v=1-\frac{526}{625}\epsilon\uparrow1,
 \qquad g\to0,\qquad v_*=4.
\]
The optimum is attained by equal mass on the latter two rows; the contrast $\theta=0$ bounds every value by four. This sequence is explicitly \emph{not} a trajectory or a counterexample to either algorithm. It only rules out a proof that substitutes bounded fits, monotone values and vanishing weighted stationarity for a boundary argument.

\section{The two sufficient steps are not universally ordered}
This exact comparison uses the same observable $G$ and compares only the selected steps, not actual improvement or total trajectory cost. Fix the intercept nodes and cache
\[
 x=(1,2,3,7/3),\quad s=(1/6,1/6,1/6,1/2),\quad
 w_\epsilon=((1-\epsilon)/3,(1-\epsilon)/3,(1-\epsilon)/3,\epsilon),
 \quad0<\epsilon<1/2.
\]
Direct substitution gives $\theta_s=-3/7$, $p=(4,1,-2,0)/7$, $A_s=91/18$, and
\[
 r=Q=0,\quad q=(1-\epsilon)/7,\quad h=2/7,\quad
 H=72/4459,\quad\alpha=2\epsilon.
\]
Therefore
\[
 G=(6+7\epsilon)/49,\quad K=36/(4459\epsilon),\quad
 B_{\rm sur}=\frac{456+588\epsilon+343\epsilon^2}{2401(1-\epsilon)}.
\]
The selected uncapped safe steps satisfy
\[
 t_K\sim(91/24)\epsilon\to0,\qquad
 t_{\rm sur}\to49/304,\qquad
 t_{\rm sur}/t_K\sim21/(494\epsilon)\to\infty.
\]
The geometry and cache are fixed; $A_{w_\epsilon}\to14/3$, so the nuisance matrix is not becoming singular. The minimum coordinate-ratio bound is responsible for this conservatism. A zero cached residual is permitted by the theorem, and every damped design stays positive.
Conversely, the earlier $(1,3,8)$ example has $K=0$, $G=3/25$, but $B_{\rm sur}=153/1250$, hence $t_K=1/2>25/102=t_{\rm sur}$. Thus neither step dominates everywhere. These are elementary coefficient comparisons, not a new optimization advantage.

\section{Full-trajectory experiment protocol and costs}
The new experiment implements three initially specified methods: uncapped relative-margin updates, fresh half-power updates, and fresh half-affine updates. Eighteen core instances use fixed features or seed 2026092202, with nuisance dimensions one through three. Nine further instances use rows $(\sqrt{c_0},0),(1,1),(1,K)$ with $K=2,10,100$. Set $c_0$ at fractions $0.001,0.5,0.999$ between the hidden two-site uniform value $(K-1)^2/[2(1+K^2)]$ and optimum $(K-1)^2/(K+1)^2$. This probes a competing face; every design begins uniformly. Each run has a 2000-update limit, with floating numerical stops retained explicitly.

The new capped $K$ policy and capped surrogate policy were subsequently tested on the saturated $(1,2)$ case, the positive cubic case, and the $K=100$, fraction-$0.999$ boundary case. The compared policies use identical features and initialization. The surrogate evaluates $R_2$ with an additional charged scalar reduction. The cap counter is the number of cache refreshes after initialization, not the iteration count. No weight floor or unreported correction changes the update.

For every instance an offline LP minimizes $\|b\|_1$ subject to $F^\top b=e_1$, giving $v_*=1/\|b\|_1^2$. Its primal/dual residuals and returned design are recorded. After every accepted update, offline objective evaluation pays for its own current matrix assembly, Cholesky factorization, solve, feature matvecs and residual-norm evaluation. The algorithm cannot access these values. First-hit and matched-resource summaries are extracted afterward; no implementation stops using the known LP optimum. The target $0.999999$ in the main table is therefore a common \emph{offline evaluation} criterion, not a free online stopping certificate.

A fresh gate exploits its known zero normal residual, using one old-factor solve for $H$; a stale gate pays for both $Q,H$. Rejected stale gates and fresh retries are both charged. Matrix assemblies, factorizations, fresh-contrast solves, cached-factor forward/back pairs, feature matvec passes and conceptual fusible scalar reductions are separate resources. The experiment makes no wall-time claim or conversion to a common FLOP scalar. Matching feature passes leaves factor and solve budgets unequal; matching those other resources is also recorded separately.

All eighteen core uncapped runs reached the target, but the fresh half-power baseline had higher efficiency at each matched feature-pass budget. Six boundary cases reached the target within 409 uncapped updates; the three near-tie cases did not by the cap. Their final hidden mass stayed approximately $0.64$--$0.68$, offering no finite evidence of a hidden-mass extinction trap. The capped cubic run reached only $0.9999101845$ at 2000 updates; its classical-surrogate counterpart reached $0.9999093154$. The latter expended 6189 scalar reductions versus 4126 for the projected coefficient, in addition to the equal counts displayed in the main text. On the difficult boundary case, the surrogate used 4060 solves and 4091 feature passes versus 4086 and 4130 for $K$, but 6090 scalar reductions versus 4086. These are resource tradeoffs and finite results, not a universal ordering.

The exact scripts additionally check twelve rational uncapped-margin steps on three models, plus two saturated half-power steps, paying for 26 exact current-value evaluations. A separate 107-assertion rational check verifies finite-interpolant identities, fresh projection bounds and the nontrajectory boundary witness: six $2\times2$ basis solves, five current $2\times2$ solves, five additional $H$ inverse actions, and five scalar normal-equation solves. Neither finite verification substitutes for the general proofs.

The machine-readable results retain numerical stops, minimum observed increments and all audit costs. Roundoff-level negative increments are not a validated certificate. A separate 49-model exploratory scan for uncapped boundary attraction found no counterexample, but is explicitly exploratory and supplies no probability or universal assurance. All executed work, including preliminary runs and the later direct-surrogate comparison, is recorded in the accompanying research cost ledger. A new three-case preconditioned-correction comparison is given below; maintained inverses, LP solve time and high-dimensional applications remain unbenchmarked.

\section{An explicit but extremely conservative accuracy--work bound}
This bound concerns precisely the capped relative-margin policy, with curvature coefficient $C=K$ or $C=B_{\rm sur}$. It is not a useful acceleration guarantee. Initialize uniformly. Divide objective-scale quantities by $v_0$; all ratios and selected steps are unchanged. Set $\ell=L/v_0$, so $1\le v\le v_*\le\ell$, $1\le q\le\ell$, $0\le\lambda_i=p_i^2\le\ell$, and $|\rho_i|\le\ell$. Fix $0<\epsilon<1$ and define a target by $v_*-v\le\epsilon$. In unnormalized units this means a gap at most $\epsilon v_0$, sufficient for efficiency at least $1-\epsilon$. The target is an analysis variable, not an online stopping input.

If $m=1$ or $\ell=1$, the initial design is optimal. Otherwise define
\begin{align*}
 a&=\epsilon/(2\ell),& B&=8\ell(\ell-1)/\epsilon^2,& A&=1+4\ell/\epsilon,\\
 J_0&=\lceil4\ell^3/\epsilon\rceil,&
 E_0&=(\tfrac12\log m+B)A^{J_0},\\
 S_{\max}&=E_0+B+(4\ell/\epsilon)(\log m+2E_0+2B),\\
 J_{\max}&=\lceil(4\ell+1)\exp(S_{\max})\rceil,&
 d_{\min}&=m^{-1}\exp(-2S_{\max}),\\
 t_{\min}&=\min\{1/2,1/(J_{\max}+1),d_{\min}^4\epsilon^2/(32\ell^4)\},\\
 N_{\max}&=\lceil S_{\max}/t_{\min}\rceil+1.
\end{align*}
The displayed arrays only define constants; $E_0$ is unrelated to the correction-baseline residual error $E$ below.

\begin{proposition}[Explicit first-hit work bound]
The first design of efficiency at least $1-\epsilon$ occurs by $N_{\max}$ accepted steps and $J_{\max}$ refreshes. Up to that hit, at most $J_{\max}+1$ caches are constructed and at most $N_{\max}+J_{\max}+1$ gates are evaluated, including rejected stale gates and fresh retries. The assertion holds for both $K$ and $B_{\rm sur}$ without a minimum-weight or matrix-conditioning hypothesis.
\end{proposition}
\begin{proof}
We bound every finite prefix whose update origins have $v_*-v>\epsilon$, allowing its last endpoint to be the first hit. The earlier exact cache satisfies $\lambda_{\min}\le v(s)\le v(w)$. Hence
\[
 \zeta=1-\lambda_{\min}/q\ge D/q,\qquad
 g/2\ge\zeta Q/\alpha\ge\zeta D\ge D^2/q.
\]
The accepted gain is at least $tg/4$, so $\sum tg\le4(\ell-1)$. Declare a step bad when $D>\epsilon/2$. Then $g>\epsilon^2/(2\ell)$ on bad steps, and their total time satisfies $\sum_{\rm bad}t\le B$.

Let $S_b=\sum_{k<b}t_k$ denote elapsed step time. Since $w_i^+\ge(1-t)w_i$ and $\log(1-t)\ge-2t$,
\[
 w_{b,i}\ge m^{-1}\exp(-2S_b).
\]
Fix an epoch beginning at $b$ and a maximum-score coordinate with $M=\max_i\lambda_i$. It never decreases because $q\le M$. At a good, subtarget-origin step, $q=v+D<v_*-\epsilon/2$ and $M\ge v_*$. Therefore
\[
 \log(w_i^+/w_i)
 \ge\frac{t(M-q)}{q+t(M-q)}\ge\frac{\epsilon t}{2\ell}=at.
\]
Here $\log(1+x)\ge x/(1+x)$ for $x\ge0$, and the denominator is at most $M\le\ell$. Bad steps have nonnegative change for this same coordinate. Its weight cannot exceed one, so the epoch's good time is at most $(\log m+2S_b)/a$. If $B_b$ is its bad time, its final clock obeys
\[
 S_{\rm end}\le A S_b+\frac{\log m}{a}+B_b.
\]
Iterating over $r$ epochs with total $\sum B_b\le B$ gives
\[
 S_{\rm end}\le\tfrac12\log m(A^r-1)+BA^{r-1}
 \le(\tfrac12\log m+B)A^r.
\]
The same bound holds for a prefix ending inside an epoch. Thus the first $J_0$ epochs have clock at most $E_0$ before the target.

For any optimal design $u_*$, let $\Phi=\mathrm{KL}(u_*\Vert w)$. Uniform initialization gives $\Phi_0\le\log m$. Every step satisfies $\Delta\Phi\le-\log(1-t)\le2t$. At the start of epoch $J_0$, if reached before the target, $\Phi\le\log m+2E_0$. On each later good step, $u_*\cdot\rho\ge\epsilon/(2\ell)=a$, so the logarithm inequality already proved yields
\[
 \Delta\Phi\le-at+\ell^2t^2\le-at/2,
\]
where the cap gives $t\le1/(J_0+1)\le\epsilon/(4\ell^3)$. Late bad steps increase entropy by at most $2B$ in total. Nonnegativity of entropy bounds late good time by
$(2/a)(\log m+2E_0+2B)$.
Adding early time and late bad time yields total clock at most $S_{\max}$ for any subtarget-origin prefix. This argument uses no unknown value-convergence time.

The accepted fresh step at refresh counter $j$ is at least
$\min\{1/(4\ell),1/(j+1)\}\ge1/(4\ell+j+1)$.
If the prefix includes the accepted fresh step at counter $J$, then
\[
 S_{\max}\ge\sum_{j=0}^J\frac1{4\ell+j+1}
 \ge\log\frac{4\ell+J+2}{4\ell+1}.
\]
This proves the conservative refresh bound $J_{\max}$. Exact optimal termination already meets the target and cannot insert an uncharged sequence of refreshes.

Every current and earlier cached weight in such a prefix is at least $d=d_{\min}$, a consequence of the proved clock bound. Thus $\alpha\ge d$ and $\max_i w_i/s_i\le1/d$. Weighted projection onto the nuisance columns gives
\[
 Q\le\sum_i(w_i^2/s_i)p_i^2\le\ell/d,\qquad
 H\le\sum_i(w_i^2/s_i)\rho_i^2p_i^2\le\ell^3/d.
\]
Consequently $K\le4\ell^3/d^2$. Also $R_2\le\ell^3$, whence
$B_{\rm sur}\le\ell^3(1+1/d)^2\le4\ell^3/d^2$.
Before the target, the cached score range exceeds $\epsilon$, because its maximum is at least $v_*$ and its minimum is at most $v(s)\le v(w)<v_*-\epsilon$. The corresponding two current weights are each at least $d$. The pairwise variance formula yields
\[
 \operatorname{Var}_w(\lambda)\ge d^2\epsilon^2,
 \quad g\ge d^2\epsilon^2/\ell,
 \quad G/(4C)\ge d^4\epsilon^2/(32\ell^4)
\]
when $C>0$. Omitting the zero-$C$ term can only increase the step. Therefore every subtarget-origin accepted step is at least $t_{\min}$. A prefix of $N$ such steps has $Nt_{\min}\le S_{\max}$; the stated $N_{\max}$ covers the first hit. An infinite subtarget run is impossible. Every refresh has at most one preceding rejected stale gate, giving the gate and cache counts.
\end{proof}
A dense arithmetic bound follows directly:
\[
 O\!\left((J_{\max}+1)(mp^2+p^3+mp)
 +(N_{\max}+J_{\max}+1)(mp+p^2+m)\right).
\]
This charges all cache builds and failed/successful gates; the surrogate has its additional scalar reduction. It is not a wall-time statement. Computing $L$ by enumerating interpolation bases, if desired, would be extra work; the actual protocol does not compute it.

For proof parameters $m=3$, $\ell=2$, $\epsilon=0.1$, one has $J_0=320$, $A=81$, $B=1600$, and $E_0>81^{320}>10^{610}$. The refresh bound then exponentiates an even larger clock bound. This example evaluates the formulas, not a design experiment. Such loose constants cannot support practical acceleration, a matching complexity lower bound, or a distinction between the two coefficients. Early-epoch control through a worst-case minimum weight and inversion of a harmonic sum are the two sources of exponential amplification.

\section{A reliable preconditioned-correction comparator}
The line-minimized inner correction is classical preconditioned steepest descent: apply Saad's Algorithm~5.2 with the symmetric preconditioning transformation in Sec.~9.2.1~\cite{saad2003} to the current normal equations, using $A_s$ as preconditioner.
At a current design, retain the original exact anchor contrast $p_0$ and old factor of $A_s$. Form $r_0=Z^\top Wp_0$. If $r_0=0$, set $p_1=p_0,r_1=0,E=0$: the fit is current-exact and no correction solve is needed. Otherwise compute
\[
 u=A_s^{-1}r_0,\quad b=Z^\top WZu,\quad
 \eta=\frac{r_0^\top u}{u^\top b}>0,\quad
 p_1=p_0-\eta Zu,\quad r_1=r_0-\eta b.
\]
The denominator is positive because $A_w\succ0$ and $u\ne0$; $\eta$ exactly minimizes the weighted quadratic along this preconditioned residual direction. Computing $b$ requires the two feature passes $Zu$ and $Z^\top WZu$, without forming $A_w$. A second old-factor solve, skipped when $r_1=0$, gives
\[
 E=r_1^\top A_s^{-1}r_1/\alpha\ge
 r_1^\top A_w^{-1}r_1.
\]
Using $p_1$, form the ordinary multiplicative direction and its $q,\rho,g,\gamma,\zeta$, then set
\[
 R_2=\sum_iw_i\rho_i^2p_{1,i}^2,\quad
 G_c=g-2\sqrt{R_2E}-\zeta E,\quad
 B_c=(\sqrt{R_2}+\gamma\sqrt E)^2.
\]
If $g>0,G_c\ge g/2$, take $t=\min\{1/2,1/(j+1),G_c/(4B_c)\}$, omitting the last term when $B_c=0$. Otherwise refresh exactly and retry. Fresh zero $g$ certifies an optimum; the fresh gate has $E=0$ and needs no correction solves. The corrected contrast is not falsely designated an exact anchor: every subsequent attempt starts from the stored original $p_0$ until an actual refresh.

For the one-step guarantee, write $e=A_w^{-1}r_1$ and $P=p_1-Ze$ for the proof-only current exact contrast. Its mismatch energy is at most $E$. With $h=Z^\top W(\rho p_1)$, projection gives $h^\top A_w^{-1}h\le R_2$. Since $A_d\succeq-\zeta A_w$,
\[
 \sum_i d_iP_i^2=g-2h^\top e+e^\top A_de\ge G_c.
\]
Weighted triangle inequality gives $\sum_iw_i\rho_i^2P_i^2\le B_c$. Substituting these two bounds into the already proved covariance-surrogate argument yields
\[
 v(w+td)-v(w)\ge tG_c/2>0.
\]
All finite accepted weights remain positive. This proof is a direct classical residual/covariance specialization. The full corrected trajectory changes its scores inside each old-factor epoch and therefore is not covered automatically by the fixed-score global or work theorem above.

A nonzero stale attempt requires at most two old-factor solves and three feature matvec passes. Compared with old-fit $K$, it uses one extra feature pass and scalar work; compared with old-fit $B_{\rm sur}$, its scalar reductions are similar but it adds nuisance-vector products for line minimization. Exact zero residuals avoid unnecessary solves. Refreshes and rejected attempts retain their costs. In scalar nuisance dimension, line minimization reaches the exact current fit; this is a property of one-dimensional quadratics, not a new solver advantage.

Only the three already selected discriminating cases were run, reusing saved offline LP values. At the positive cubic target $0.999999$, correction required 879 updates and 2658 feature passes; fresh half-power required 203 and 406. At the old capped-$K$ budget of 4190 feature passes, correction efficiency was $0.9999999852$, compared with $0.9999101845$ for $K$ and $0.9999093154$ for the uncorrected surrogate. At the difficult boundary's 4130-pass budget, correction gave $0.9997838744$, versus $0.9997690537$ and $0.9997717331$. These finite comparisons match one resource only and preserve all other costs in the evidence files. They do not establish superiority over the fresh methods or globally improved complexity.

\section{Quantitative comparisons with established methods}
In this comparison, $N$ is the candidate count and $r$ the full regression dimension. Richt\'arik~\cite{richtarik2008}, Algorithm~2 and Theorem~34 (report PDF pages 34--36), assumes nonzero spanning candidate vectors, a nonzero target and a uniform start. With $\psi(w)=\sqrt{c^\top M(w)^{-1}c}$, its guarantee is
\[
 \psi(w_{\rm out})\le(1+\delta)\psi_*,\qquad
 K_{\rm it}\le2\Gamma\,[\log\Gamma+\log\log N+8/\delta].
\]
The selected full-matrix leverage must satisfy
$f_{j_k}^\top M(w_k)^{-1}f_{j_k}\le\Gamma$ throughout the run (Assumption~33). This is not implied by bounded target residuals: nuisance eigenvalues can approach zero. Since $v=\psi^{-2}$, the resulting efficiency is at least $(1+\delta)^{-2}$; the target $1-\epsilon$ uses $\delta=(1-\epsilon)^{-1/2}-1$. The report's Theorem~38 instead concerns the objective bias from a permanent weight floor; it does not remove Assumption~33 for the original feasible problem.

The fresh coordinate update has
\[
 w^+=\frac{w+\kappa e_j}{1+\kappa},\qquad
 M^+=\frac{M+\kappa f_jf_j^\top}{1+\kappa}.
\]
Dense factor maintenance and target-solution updating cost $O(r^2)$ per rank-one step, and a full current-sensitivity scan costs $O(Nr)$. Including initialization gives the explicit implementation accounting
$O(Nr^2+r^3+K_{\rm it}(Nr+r^2))$.
This arithmetic count follows from the displayed updates, not from treating a fresh factorization as free or asserting numerical stability. An infinite line parameter in the original algorithm is an optimal terminal branch, not a requirement to factor a singular limiting matrix and continue. One-coordinate updates and simultaneous reweighting have different iteration counts. Their common-accuracy total work still needs measurement.

For John ellipsoids, let $A\in\R^{N\times r}$ have full column rank and $\epsilon,\delta\in(0,1)$. Cohen et al.~\cite{cohen2019}, revised preprint Theorem~3.6/C.7, give $O(\epsilon^{-1}\log(N/\delta))$ iterations with $O(\epsilon^{-1})$ current linear systems per iteration and probability at least $1-\delta$. The output is averaged and normalized, with total weight $r$ and maximum leverage at most $1+\epsilon$. Converting it to simplex weights gives maximum leverage at most $r(1+\epsilon)$; this does not imply $1-\epsilon$ efficiency for an arbitrary fixed target $c$. System counts also require solver and matrix-application costs to become arithmetic counts. For $P=\{x:\|Ax\|_\infty\le1\}$, Woodruff and Yasuda~\cite{woodruff2024}, Theorem~1.6, return an ellipsoid $\mathcal E$ with $(1+\epsilon)^{-1/2}\mathcal E\subseteq P\subseteq\sqrt r\,\mathcal E$ with probability at least $1-1/\operatorname{poly}(N)$. Their lazy total-runtime bound involves all of
\[
\begin{split}
 &O(\epsilon^{-1}Nr)[\log(N/r)+\operatorname{polylog}(\epsilon^{-1}\log N)]\\
 &\quad+O(N)\operatorname{poly}(\epsilon^{-1}\log N)
 +O(N^{0.1})r^{\omega+1}\epsilon^{-3}(\log N)^2.
\end{split}
\]
Its randomized sketches, matrix multiplication and batched recovery do not describe an unchanged cached $c$-contrast. Dropping the last two terms or equating the two accuracy objectives would misstate this comparison. No unsupported last-iterate conclusion is imported: the cited John algorithms use averaged output.

There is also a standard quantitative use of the present old factor within one fixed current design. This is preconditioned Richardson iteration; see Saad~\cite[Example~4.1, Eq.~(4.33), and Sec.~9.2.1]{saad2003}. Here $[\alpha,\beta]$ is a certified spectral enclosure, so the displayed factor is a valid upper bound without computing the extreme eigenvalues. Since $\alpha A_s\preceq A_w\preceq\beta A_s$, the stationary correction
\[
 \theta^{(k+1)}=\theta^{(k)}-\omega A_s^{-1}r(\theta^{(k)}),\qquad
 \omega=\frac{2}{\alpha+\beta},\quad
 \chi=\frac{\beta-\alpha}{\beta+\alpha}
\]
satisfies, from the symmetric error operator
$I-\omega A_w^{1/2}A_s^{-1}A_w^{1/2}$,
\[
 q_w(\theta^{(k)})-v(w)
 \le\chi^{2k}[q_w(\theta^{(0)})-v(w)]
 \le\chi^{2k}Q/\alpha.
\]
Each correction uses one old-factor solve and $O(mp)$ feature work. For $0<\chi<1$ and $Q/\alpha>\xi>0$, a sufficient number is
$\lceil\log((Q/\alpha)/\xi)/(-2\log\chi)\rceil$.
If the initial bound is already below $\xi$, no correction is needed; if $\alpha=\beta$, one correction is exact. This elementary preconditioned spectral calculation is background, not a separate design contribution. The line-minimized one-correction comparison above spends at most two old solves including its residual certificate. Neither calculation proves that these inner corrections remain inexpensive along every design trajectory.

\section{An actual infinite-refresh counterexample}\label{app:trap}
This section proves the main counterexample for the original uncapped $K$ protocol. All logarithms are natural and terminating decimals in bounds are exact rational numbers. The finite multiplicities are specified analytically; no enormous feature matrix or trajectory is claimed to have been numerically executed. Labels (T1)--(T56) refer only to this construction. The proof follows every actual acceptance, variable step and refresh.
\subsection{The feature matrix}

Set

\[
R=100,\qquad B=\frac{101}{100},\qquad \delta=\frac{99}{100},
\]

and

\[
N=2^{10^7},\qquad m=N^3+N^2+N+1.
\]

Take nuisance dimension \(p=1\) and the following rows:

\begin{center}
\begin{tabular}{lll}\toprule
Type & Feature row & Multiplicity \\ \midrule
Base & \((1,0)\) & \(N^3\) \\
Auxiliary & \((\sqrt{\delta},0)\) & \(N^2\) \\
First informative row & \((\sqrt B,-R)\) & \(1\) \\
Second informative row & \((\sqrt B,1)\) & \(N\) \\
\bottomrule\end{tabular}
\end{center}

Initialize every individual row with weight \(1/m\), and cache exactly there.

The feature matrix has rank \(2\). Write the four \textbf{aggregate} weights as

\[
(w_0,\omega,x_1,x_2).
\]

Their initial values are

\[
(w_0,\omega,x_1,x_2)
=\frac1m(N^3,N^2,1,N).
\tag{T1}
\]

\subsubsection{Aggregation preserves the exact protocol}

Copies of the same row always have the same weight. If a row type has \(n_j\) copies and aggregate current/cache weights \(W_j,S_j\), then each copy has weights \(W_j/n_j,S_j/n_j\), so

\[
\frac{w_i}{s_i}=\frac{W_j}{S_j},
\qquad
\frac{d_i}{w_i}=\frac{p_j^2}{q}-1.
\]

All sums defining \(A_s,q,r,h,Q,H,g\) aggregate exactly, and so do the extrema defining \(\alpha,\gamma,\zeta\). Therefore the four-weight analysis below is exactly the specified algorithm on the \(m\)-row matrix.

For positive aggregate weights,

\[
A_w=R^2x_1+x_2>0,
\]

and direct elimination gives

\[
v(w)
=w_0+\delta\omega+
B(R+1)^2\frac{x_1x_2}{R^2x_1+x_2}.
\tag{T2}
\]

The optimum is

\[
\boxed{v^*=B.}
\tag{T3}
\]

Indeed, evaluation at \(\theta=0\) gives \(v(w)\le B\) for every design. Conversely, the design supported on the informative pair with

\[
x_1=\frac1{R+1},\qquad x_2=\frac R{R+1}
\]

has minimizing coefficient \(\theta=0\) and value \(B\).

Also, the stopping test never fires: the two positive-weight, zero-nuisance row types always have distinct cached squared residuals \(1\) and \(\delta\). Thus \(g>0\) at every iterate, including every refresh.



\subsection{An epoch lemma for the actual rule}

An epoch starts with a fresh cache and ends at the first state where that old cache fails \(G\ge g/2\), before refreshing there.

During an epoch, label the two informative rows as follows:

\(x\) is the weight of the row that is initially much smaller;\par
\(y\) is the weight of the other informative row.\par

The labels remain fixed within the epoch. There are two types of epoch: \(\mathrm H\), when \(x=x_1\), and \(\mathrm L\), when \(x=x_2\).

Define

\[
\mathcal A_{\mathrm H}=B(R+1)^2=\frac{1030301}{100},
\qquad
\mathcal A_{\mathrm L}=B(1+1/R)^2=\frac{1030301}{10^6},
\tag{T4}
\]

and

\[
a_\sigma=\log\frac{1+\mathcal A_\sigma}{1+\delta},
\qquad
\beta=\log(1+\delta),
\qquad
C_\sigma=\frac{a_\sigma+2\beta}{2a_\sigma+\beta}.
\tag{T5}
\]

\subsubsection{Epoch lemma}

Suppose a fresh epoch starts with

\[
\omega_s\le10^{-6},\qquad
-\log\frac{s_y}{\omega_s}=b\ge1000,
\qquad
-\log\frac{s_x}{\omega_s}=2b+\varepsilon,
\quad |\varepsilon|\le100.
\tag{T6}
\]

Then this epoch ends after finitely many steps. Throughout the epoch, including its ending state,

\[
\omega\le10^{-6},
\qquad
\frac{x}{\omega},\frac{y}{\omega}\le10^{-20}.
\tag{T7}
\]

The auxiliary weight \(\omega\) decreases, the first step is exactly \(1/2\), and every accepted step is at least \(10^{-12}\).

At the ending state, put

\[
b'=-\log\frac{x_{\rm end}}{\omega_{\rm end}}.
\]

Then

\[
\left|
-\log\frac{y_{\rm end}}{\omega_{\rm end}}-2b'
\right|\le100,
\qquad
|b'-C_\sigma b|\le1000.
\tag{T8}
\]

Consequently, whenever \(b'\ge1000\), interchanging the labels of the informative rows gives the hypotheses for the next epoch.

Here is a proof with explicit estimates.

\subsubsection{2.1. Cached residuals}

Write

\[
\rho=\frac{s_x}{s_y}=e^{-b-\varepsilon},
\qquad
\kappa=\frac{z_x^2}{z_y^2}\in\{10^4,10^{-4}\}.
\]

The two informative rows have opposite nuisance signs. Solving the scalar normal equation gives their cached squared residuals

\[
P:=p_x^2=\frac{\mathcal A_\sigma}{(1+\kappa\rho)^2},
\qquad
S:=p_y^2=\kappa\rho^2P.
\tag{T9}
\]

Moreover,

\[
A_s=z_y^2s_y(1+\kappa\rho).
\tag{T10}
\]

Since \(b\ge1000\) and \(|\varepsilon|\le100\),

\[
1.02<P<11000,\qquad 0<S<10^{-20},
\qquad 1\le1+\kappa\rho\le2.
\tag{T11}
\]

Thus the cached squared sensitivities are strictly ordered:

\[
S<\delta<1<P.
\tag{T12}
\]

\subsubsection{2.2. Estimates in the region (T7)}

First derive estimates conditional on (T7); an induction below establishes that the trajectory remains there.

Because \(w_0=1-\omega-x-y\),

\[
q=1-(1-\delta)\omega+(P-1)x+(S-1)y.
\]

Consequently,

\[
1-.011\omega\le q\le1-.009\omega,
\qquad .999<q<1.
\tag{T13}
\]

The identity

\[
g=\frac1q\sum_i w_i(p_i^2-q)^2
\]

then gives

\[
10^{-5}\omega\le g\le10^{-3}\omega.
\tag{T14}
\]

For the lower bound, the auxiliary row alone contributes
\(\omega(q-\delta)^2/q\), and \(q-\delta>.009\).
For the upper bound, the auxiliary contribution is below
\(0.000101\omega\), the base-row contribution is at most
\(.011^2\omega^2/.999\), and the two informative contributions are bounded using \(x,y\le10^{-20}\omega\) and \(P<11000\).

For a step \(t\), write

\[
F_U=1+t(U/q-1)
\]

for the multiplier of a row with cached squared residual \(U\). For \(0<t\le1/2\), (T11)–(T13) imply

\[
\frac{F_P}{F_\delta}<10^4,\qquad
\frac{F_S}{F_\delta}<1,\qquad
\frac{F_P^2}{F_\delta F_S}<10^{10},
\tag{T15}
\]

and

\[
F_\delta\le1-.009t.
\tag{T16}
\]

For example, \(F_P<5600\), \(F_\delta\ge.995\), and \(F_S\ge1/2\).
When \(t=1/2\), (T16) gives \(F_\delta<.996\).

Introduce the positive balance variable

\[
W=\frac{x^2}{\omega y}.
\tag{T17}
\]

\subsubsection{2.3. The first step really is \texorpdfstring{$1/2$}{one half}}

At the fresh state,

\[
W_0=\frac{s_x^2}{\omega_s s_y}
=e^{-3b-2\varepsilon}\le e^{-2800}.
\tag{T18}
\]

The cache normal equation gives

\[
z_xp_xs_x=-z_yp_ys_y.
\]

Hence, at the fresh state,

\[
H=
\frac{\kappa P}{1+\kappa\rho}
\frac{s_x^2}{s_y}
\left(\frac{P-S}{q}\right)^2.
\]

Using (T11), (T13), and (T14),

\[
\frac Hg\le 2\cdot10^{22}W_0<\frac12.
\tag{T19}
\]

Since fresh \(Q=0\), \(G=g\), and \(K=H\), the prescribed step is exactly

\[
t=\frac12.
\tag{T20}
\]

After this step, the relative current-to-cache growth of the two informative rows satisfies

\[
J:=\frac{x/s_x}{y/s_y}
=\frac{q+P}{q+S}>2.
\tag{T21}
\]

It increases strictly at all subsequent steps in this epoch.

\subsubsection{2.4. Exact control of \texorpdfstring{$\alpha,\gamma,\zeta,G,K$}{alpha, gamma, zeta, G, K}}

Let

\[
\eta_x=x/s_x,\qquad \eta_y=y/s_y.
\]

Because of the fixed ordering (T12), the products of update multipliers preserve the same ordering of current-to-cache ratios. In particular,

\[
\boxed{\alpha=\eta_y.}
\tag{T22}
\]

This includes the base and auxiliary rows and all their copies.

Let

\[
\mathcal C=z_xp_xs_x=-z_yp_ys_y,
\qquad
\Delta_x=P/q-1,\quad \Delta_y=S/q-1.
\]

At every nonfresh state,

\[
r=\mathcal C(\eta_x-\eta_y),
\qquad
h=\mathcal C(\eta_x\Delta_x-\eta_y\Delta_y).
\]

Therefore, using \(J>2\),

\[
\ell:=\frac hr
=\Delta_x+\frac{\Delta_x-\Delta_y}{J-1}
\quad\text{satisfies}\quad
0<\ell<10^5.
\tag{T23}
\]

Thus \(\sqrt H=\ell\sqrt Q\). Also,

\[
\gamma<10^5,\qquad
\zeta=1-S/q\in[1/2,1].
\tag{T24}
\]

A direct substitution into \(Q=r^2/A_s\), using (T10) and (T22), gives the exact formula

\[
\frac Q\alpha
=
\frac{\kappa P}{1+\kappa\rho}
(1-J^{-1})^2\frac{x^2}{y}.
\tag{T25}
\]

The coefficient multiplying \(x^2/y\) lies between \(10^{-6}\) and \(10^9\). Combining this with (T14), and defining

\[
\mathcal Z=\frac{Q}{\alpha g},
\]

yields

\[
10^{-3}W\le\mathcal Z\le10^{14}W.
\tag{T26}
\]

Crucially, the actual certificate quantities now satisfy

\[
1-\frac Gg=(2\ell+\zeta)\mathcal Z,
\qquad
\frac Kg=(\ell+\gamma)^2\mathcal Z.
\tag{T27}
\]

In particular,

\[
1-\frac Gg\le10^6\mathcal Z,
\qquad
\frac Kg\le4\cdot10^{10}\mathcal Z.
\tag{T28}
\]

At any accepted nonfresh state, \(G/g\ge1/2\). Since
\(2\ell+\zeta\ge1/2\), (T27) implies \(\mathcal Z\le1\). Therefore

\[
W\le10^3
\tag{T29}
\]

and the prescribed step satisfies

\[
t=\min\left(\frac12,\frac{G}{4K}\right)
\ge
\min\left(\frac12,\frac1{32\cdot10^{10}}\right)
>10^{-12}.
\tag{T30}
\]

No lower bound on any weight was used.

\subsubsection{2.5. Invariance of the region and finiteness of the epoch}

The starting state satisfies (T7). Its fresh half-step remains in that region by (T15), (T18).

For the subsequent induction, \(y/\omega\) decreases by (T15), so

\[
y/\omega\le e^{-b}.
\]

At every accepted nonfresh state, (T29) gives

\[
\frac{x}{\omega}
=\sqrt{W\,\frac y\omega}
\le\sqrt{10^3}\,e^{-b/2}.
\]

After the next step, (T15) therefore gives

\[
\frac{x_{\rm next}}{\omega_{\rm next}}
\le10^4\sqrt{10^3}\,e^{-b/2}
<10^{-20}.
\]

The other ratio also remains below \(10^{-20}\), and \(\omega\) decreases by (T16). This proves (T7) inductively, including the first failed state.

To prove that failure occurs after finitely many steps, define

\[
a_k=\log(F_P/F_\delta)>0,
\qquad
\beta_k=\log(F_\delta/F_S)>0.
\]

Then

\[
\log W_{k+1}-\log W_k=2a_k+\beta_k.
\tag{T31}
\]

Moreover,

\[
F_\delta-F_S=\frac{t(\delta-S)}q>.98t,\qquad F_S\le1,
\]

so

\[
\beta_k\ge\log(1+.98t)\ge\frac t2.
\tag{T32}
\]

Every accepted step thus increases \(\log W\) by at least
\(5\cdot10^{-13}\), whereas every accepted state has \(W\le10^3\), including the fresh state. Infinite retention of this cache is impossible.

For completeness, an explicit upper bound on the number \(M\) of steps in the epoch is

\[
M\le
2+\left\lceil
2\cdot10^{12}\bigl(3b+200+\log(10^3)\bigr)
\right\rceil.
\tag{T33}
\]

Thus no unknown hitting time is being assumed.

At the ending state, \(g>0\) and \(G/g<1/2\). Equations (T26) and (T28) give

\[
\frac12<10^6\mathcal Z\le10^{20}W_{\rm end},
\]

so \(W_{\rm end}>10^{-22}\). The preceding accepted state had \(W\le10^3\), and the one-step multiplier in (T15) is below \(10^{10}\). Hence

\[
10^{-22}<W_{\rm end}\le10^{13}.
\tag{T34}
\]

In particular,

\[
|\log W_{\rm end}|\le100.
\tag{T35}
\]

\subsubsection{2.6. The logarithmic epoch map, with a uniform error bound}

It remains to prove the second estimate in (T8). This requires controlling the actual variable steps, not replacing them by half-steps.

Whenever a nonfresh state has \(W\le10^{-30}\), (T26)–(T28) give

\[
\mathcal Z\le10^{-16},\qquad
G/g\ge1-10^{-10},\qquad
K/g\le4\cdot10^{-6}.
\]

Therefore that state is accepted and its prescribed step is exactly \(1/2\).

Let \(n\) be the first index with \(W_n>10^{-30}\). The first \(n\) steps are consequently all half-steps. From (T15),

\[
W_n\le10^{-20}<1.
\]

Each such step increases \(\log W\) by at least \(1/4\), by (T32). Since
\(-\log W_0=3b+2\varepsilon\),

\[
n\le20b.
\tag{T36}
\]

The remaining part of the epoch has uniformly bounded logarithmic motion:

\[
\sum_{k=n}^{M-1}(2a_k+\beta_k)
=\log(W_{\rm end}/W_n)
<43\log10<150.
\tag{T37}
\]

Thus the tail contributes less than \(75\) to \(\sum a_k\) and less than \(150\) to \(\sum\beta_k\), regardless of its number or sizes of steps.

For the half-step prefix, the exact increments are

\[
a_k=\log(q_k+P)-\log(q_k+\delta),
\qquad
\beta_k=\log(q_k+\delta)-\log(q_k+S).
\tag{T38}
\]

The corresponding constants at \(q=1,P=\mathcal A_\sigma,S=0\) are precisely \(a_\sigma,\beta\) from (T5).

From (T9),

\[
|P-\mathcal A_\sigma|
\le2.2\cdot10^8e^{-b+100},
\qquad
S\le1.1\cdot10^8e^{-2b+200}.
\tag{T39}
\]

Also, throughout the half-step prefix,
\(\omega_{k+1}<.996\omega_k\), so (T13) gives

\[
\sum_{k<n}|q_k-1|
\le .011\sum_{k<n}\omega_k
\le2.75\omega_s
\le2.75\cdot10^{-6}.
\tag{T40}
\]

The relevant logarithmic derivatives in (T38) have absolute value at most \(2\). Equations (T36), (T39), and (T40), with \(b\ge1000\), therefore give

\[
\left|\sum_{k<n}a_k-na_\sigma\right|<10,
\qquad
\left|\sum_{k<n}\beta_k-n\beta\right|<10.
\tag{T41}
\]

These are explicit uniform bounds. The two residual-error sums are at most
\[
 E_P=20b(2.2\cdot10^8)e^{-b+100}<1,\qquad
 E_S=20b(1.1\cdot10^8)e^{-2b+200}<1.
\]
Both functions decrease for $b\ge1000$. At $b=1000$, replacing $e^{-900}$ and $e^{-1800}$ by the larger $2^{-900}$ and $2^{-1800}$ already gives the strict bounds. Each deviation $|a_k-a_\sigma|$ or $|\beta_k-\beta|$ from the corresponding limiting increment in (T38) is bounded by $4|q_k-1|$ plus twice its relevant residual error; hence each sum in (T41) is at most $4(2.75\cdot10^{-6})+2\max(E_P,E_S)<10$. This controls the entire half-step prefix uniformly in its length.

Put

\[
X=-\log(x/\omega),\qquad Y=-\log(y/\omega).
\]

Combining the prefix estimates with (T37) gives

\[
X_{\rm end}=2b+\varepsilon-na_\sigma+E_X,
\qquad
Y_{\rm end}=b+n\beta+E_Y,
\qquad
|E_X|,|E_Y|\le200.
\tag{T42}
\]

By (T17) and (T35),

\[
\varepsilon_{\rm end}:=Y_{\rm end}-2X_{\rm end}
=\log W_{\rm end},
\qquad |\varepsilon_{\rm end}|\le100.
\tag{T43}
\]

Eliminating \(n\) from (T42) yields

\[
\begin{aligned}
X_{\rm end}
={}&C_\sigma b
+\frac{\beta}{2a_\sigma+\beta}\varepsilon
-\frac{a_\sigma}{2a_\sigma+\beta}\varepsilon_{\rm end}\\
&+\frac{a_\sigma}{2a_\sigma+\beta}E_Y
+\frac{\beta}{2a_\sigma+\beta}E_X.
\end{aligned}
\tag{T44}
\]

Every displayed coefficient of an error has absolute value below \(1\). Thus

\[
|X_{\rm end}-C_\sigma b|\le1000.
\]

Together with (T43), this proves (T8) and completes the epoch lemma.



\subsection{The two-epoch map expands}

For the chosen constants,

\[
a_{\mathrm H}=\log\frac{1030401}{199},
\qquad
a_{\mathrm L}=\log\frac{2030301}{1990000},
\qquad
\beta=\log\frac{199}{100}.
\tag{T45}
\]

We need only the elementary bounds

\[
0<a_{\mathrm H}<9,\qquad
0<a_{\mathrm L}<.021,\qquad
\beta>.68.
\tag{T46}
\]

For the first, \(1030401/199<5200<e^9\).
For the second, use \(\log(1+x)<x\).
For the third, putting \(z=99/299\),

\[
\log\frac{199}{100}
=2\operatorname{arctanh}z
>2\left(z+\frac{z^3}{3}\right)
=\frac{18348264}{26730899}>.68.
\]

The function

\[
C(a,\beta)=\frac{a+2\beta}{2a+\beta}
\]

decreases in \(a>0\) and increases in \(\beta>0\). Therefore

\[
\begin{aligned}
C_{\mathrm H}C_{\mathrm L}
&>
\frac{9+1.36}{18+.68}
\frac{.021+1.36}{.042+.68}\\
&=\frac{357679}{337174}
>\frac{21}{20}=1.05.
\end{aligned}
\tag{T47}
\]

Also,

\[
\frac12<C_\sigma<2.
\tag{T48}
\]

Suppose a high-type epoch starts with parameter \(b\ge10^6\). The next, low-type epoch starts with parameter

\[
b'\ge C_{\mathrm H}b-1000
>\frac b2-1000\ge499000,
\tag{T49}
\]

so the epoch lemma remains applicable. The following high-type epoch starts with

\[
\begin{aligned}
b''
&\ge C_{\mathrm L}b'-1000\\
&\ge C_{\mathrm L}C_{\mathrm H}b
-1000C_{\mathrm L}-1000\\
&>1.05b-3000
\ge1.04b.
\end{aligned}
\tag{T50}
\]

Thus the high-type epoch parameters grow geometrically. Every intervening low-type epoch also remains within the lemma’s hypotheses. At each failed endpoint the actual gate provides a new error parameter with absolute value at most $100$; this parameter is reset by the endpoint identity and is not an error accumulated over all preceding epochs. Moreover $2b'+\varepsilon_{\rm next}>b'$ when $b'\ge1000$, so the two informative rows really exchange their small/large roles at every refresh.

This proves an infinite succession of finite epochs, alternating between the two informative rows.



\subsection{The uniform initialization enters this induction exactly}

For (T1),

\[
\frac{x_1}{\omega}=N^{-2},
\qquad
\frac{x_2}{\omega}=N^{-1}.
\]

The initial epoch is high-type and satisfies

\[
X_0=2\log N,\qquad Y_0=\log N.
\]

Thus its parameters are

\[
b_0=\log N=10^7\log2>5\cdot10^6,
\qquad \varepsilon_0=0.
\tag{T51}
\]

Moreover,

\[
\omega_0=\frac{N^2}{m}<\frac1N<10^{-6}.
\tag{T52}
\]

Equations (T51)–(T52) satisfy the hypotheses above. Consequently, the \textbf{exact uniformly initialized algorithm} has infinitely many refreshes, and the region

\[
x_1/\omega,\ x_2/\omega\le10^{-20}
\tag{T53}
\]

holds at every iterate.

All individual weights remain strictly positive: every update multiplier is at least \(1-t\ge1/2\). Every cache factorization is therefore well-defined.

Each epoch contains a fresh half-step, which decreases \(\omega\) by a factor below \(.996\); all other steps decrease it as well. Since there are infinitely many epochs,

\[
\omega_k\longrightarrow0.
\]

Equation (T53) then gives

\[
x_{1,k}\longrightarrow0,\qquad
x_{2,k}\longrightarrow0,\qquad
w_{0,k}\longrightarrow1.
\tag{T54}
\]

For every iterate, nonnegativity of the informative-pair contribution and evaluation at \(\theta=0\) give

\[
w_0\le v(w)
\le w_0+\delta\omega+B(x_1+x_2)
=1-.01\omega+.01(x_1+x_2)<1.
\tag{T55}
\]

Together with (T54),

\[
\boxed{\lim_{k\to\infty}v(w_k)=1<\frac{101}{100}=v^*.}
\tag{T56}
\]

The established gain bound, with \(g>0\), also shows that every accepted step strictly increases \(v\). Hence this is genuinely

\[
v(w_k)\uparrow1,
\]

not a failure of monotonicity or a finite-precision support deletion.


\section{The same boundary trap for the uncapped surrogate step}
\label{app:bsur-trap}

This section considers a second actual trajectory on the finite instance of
Section~\ref{app:trap}. The slope certificate, relative margin, and exact-refresh
semantics are unchanged. Only the curvature coefficient in the accepted step is
replaced by
\[
 R_2=\sum_iw_i\left(\frac{p_i^2}{q}-1\right)^2p_i^2,
 \qquad B_{\rm sur}=\left(\sqrt{R_2}
                  +\gamma\sqrt{Q/\alpha}\right)^2.
\]
Thus a cache is retained exactly when $g>0$ and $G\ge g/2$. Otherwise it is
recomputed at the current design; only a fresh $g=0$ terminates the run. An
accepted update uses
\begin{equation}\label{eq:bsur-rule}
 t=\min\{1/2,G/(4B_{\rm sur})\},\qquad w^+=w+td.
\end{equation}
The quotient is omitted if $B_{\rm sur}=0$, a case that will not occur in the
construction. There is no refresh-count cap. The old-cache quantities in this
rule are exactly
\[
 q=\sum_iw_ip_i^2,\quad d_i=w_i(p_i^2/q-1),\quad
 g=\frac{\operatorname{Var}_w(p_i^2)}q,\quad
 G=g-\frac{2\sqrt{HQ}+\zeta Q}{\alpha}.
\]
In particular, the cached fit is not silently replaced by a current fit when
evaluating the gate or the step.

\paragraph{Claim and fixed instance.}
Set $R=100$, $B=101/100$, $\delta=99/100$, and $N=2^{10^7}$. The rows
\[
 (1,0),\quad(\sqrt\delta,0),\quad
 (\sqrt B,-R),\quad(\sqrt B,1)
\]
have respective multiplicities $N^3,N^2,1,N$, and all
$m=N^3+N^2+N+1$ individual rows initially have weight $1/m$.
For the rule~\eqref{eq:bsur-rule}, every finite weight stays positive, every cache
epoch is finite, exact refreshes occur infinitely often, and
\begin{equation}\label{eq:bsur-claim}
 v(w_k)\uparrow1<101/100=v_*,\qquad
 \operatorname{eff}(w_k)\longrightarrow100/101.
\end{equation}
Here $k$ counts accepted updates. The assertion permits repeated rows and does
not assert uniform initialization on the four distinct types after deduplication.

Equal copies remain equal under the update. All sums, now including $R_2$, and
all extrema defining $\alpha,\gamma,\zeta$ aggregate exactly. The four aggregate
masses $(w_0,\omega,x_1,x_2)$ therefore start at
\[
 (w_0,\omega,x_1,x_2)=m^{-1}(N^3,N^2,1,N).
\]
The feature matrix has full column rank, and direct elimination gives
\begin{equation}\label{eq:bsur-value}
 A_w=R^2x_1+x_2>0,\qquad
 v(w)=w_0+\delta\omega+
 B(R+1)^2\frac{x_1x_2}{R^2x_1+x_2}.
\end{equation}
Evaluation at $\theta=0$ bounds every design value by $B$. The informative-pair
design $(x_1,x_2)=(1/(R+1),R/(R+1))$ attains this bound, so $v_*=B$.
The two zero-nuisance types have distinct fixed squared residuals $1,\delta$;
their positive weights ensure $g>0$ at every finite state, including refreshes.

\subsection{An actual Bsur epoch: geometry and finite exit}
\label{sec:bsur-geometry}

At a fresh cache, let $x$ denote the initially smaller informative mass and $y$
the other one, keeping these labels fixed throughout the epoch. Its type is
$\mathrm H$ if $x=x_1$ and $\mathrm L$ if $x=x_2$. Assume
\begin{equation}\label{eq:bsur-entry}
 0<\omega_s\le10^{-6},\qquad
 s_y/\omega_s=e^{-b},\qquad s_x/\omega_s=e^{-2b-\varepsilon},
 \qquad b\ge1000,\quad |\varepsilon|\le100.
\end{equation}
The epoch ends at the first state where the old cache fails its relative margin,
before taking a fresh fit there. We first prove invariance, a positive step
lower bound, and finite exit for the actual Bsur steps.

Write $\rho=s_x/s_y=e^{-b-\varepsilon}$,
$\kappa=z_x^2/z_y^2\in\{10^4,10^{-4}\}$, and
\[
 \mathcal A_{\mathrm H}=1030301/100,\qquad
 \mathcal A_{\mathrm L}=1030301/10^6.
\]
The fresh scalar normal equation gives the fixed cached scores
\begin{equation}\label{eq:bsur-scores}
 P=p_x^2=\frac{\mathcal A_\sigma}{(1+\kappa\rho)^2},\qquad
 S=p_y^2=\kappa\rho^2P,\qquad
 A_s=z_y^2s_y(1+\kappa\rho).
\end{equation}
In particular,
\[
 1.02<P<11000,\qquad 0<S<10^{-20},\qquad
 1\le1+\kappa\rho\le2,\qquad S<\delta<1<P.
\]
These are fresh-fit identities; they do not assert equality with the trajectory
using $K$.

\paragraph{Conditional regional bounds.}
First suppose a current state belongs to the region
\begin{equation}\label{eq:bsur-region}
 \omega\le10^{-6},\qquad x/\omega\le10^{-20},\qquad y/\omega\le10^{-20}.
\end{equation}
An induction below proves this condition through the first failed endpoint.
Since $w_0=1-\omega-x-y$,
\[
 q=1-.01\omega+(P-1)x+(S-1)y,
 \qquad 1-.011\omega\le q\le1-.009\omega,
 \qquad .999<q<1.
\]
The auxiliary term in $qg=\sum_iw_i(p_i^2-q)^2$ gives
$g\ge10^{-5}\omega$. For an upper bound, its contribution to $g$ is below
$.000101\omega$, the base contribution is at most
$.011^2\omega^2/.999$, and the other contributions total at most
$(11000^2+1)10^{-20}\omega/.999$. Thus
\begin{equation}\label{eq:bsur-g-bounds}
 10^{-5}\omega\le g<.000102\omega<10^{-3}\omega.
\end{equation}
The new curvature statistic is exactly
\[
 R_2=\frac{w_0(1-q)^2+\delta\omega(\delta-q)^2
                 +Px(P-q)^2+Sy(S-q)^2}{q^2}.
\]
Put $U=R_2/g$. The auxiliary contribution gives
$R_2>\delta(.009)^2\omega=.00008019\omega$;
since $.00008019>(3/4)(.000102)$, we have $U>3/4$.
For the upper bound use the exact centered-moment identity
\[
 q^2(R_2-g)=\sum_iw_i(\lambda_i-q)^3,
 \qquad \lambda=(1,\delta,P,S).
\]
The auxiliary contribution is less than $-(.009)^3\omega$, the base contribution
is at most $.011^3\omega^3$, the $x$ contribution is at most
$11000^3\,10^{-20}\omega$, and the $y$ contribution is negative. Hence
\[
 q^2(R_2-g)<\omega[-.009^3+.011^3\,10^{-12}+11000^3\,10^{-20}]<0.
\]
We have proved the step-independent regional bounds
\begin{equation}\label{eq:bsur-ratio-region}
 3/4<U<1,
\end{equation}
and $R_2>0$.

For $0<t\le1/2$, let $F_V=1+t(V/q-1)$ denote a row multiplier. The same
regional bounds give
\begin{equation}\label{eq:bsur-multipliers}
 \frac{F_P}{F_\delta}<10^4,\quad
 \frac{F_S}{F_\delta}<1,\quad
 \frac{F_P^2}{F_\delta F_S}<10^{10},\quad
 F_\delta\le1-.009t.
\end{equation}
For example, $F_P<5600$, $F_\delta\ge.995$, and $F_S\ge1/2$ suffice for the
two upper ratio bounds.

\paragraph{The actual fresh step and informative growth.}
At the fresh origin, $Q=0$, $G=g$, and $B_{\rm sur}=R_2$. Therefore
\begin{equation}\label{eq:bsur-fresh-step}
 t_{\rm fresh}=1/(4U)\in(1/4,1/3),\qquad
 \omega^+<.99775\omega.
\end{equation}
This step is not a half-step. After it,
\[
 J:=\frac{x/s_x}{y/s_y}
 =\frac{F_P}{F_S}
 =1+\frac{t(P-S)}{q(1-t)+tS}>4/3.
\]
Indeed, the numerator of $J-1$ exceeds $(1/4)(1.02-10^{-20})$ and its
denominator is below $3/4+(1/3)10^{-20}$. Three times that lower numerator
exceeds the upper denominator. Throughout the remaining epoch, $J$ increases.
In a low-type fresh epoch $J<1+(1/2)\mathcal A_{\rm L}/.999<1.52$;
thus the earlier $K$ proof's $J>2$ cannot be used for this rule.

\paragraph{The actual gate and step floor.}
The ordering of the fixed scores preserves the ordering of current-to-cache
ratios, including the base and auxiliary types. Thus, with
$\eta_x=x/s_x$ and $\eta_y=y/s_y$, we have $\alpha=\eta_y$.
Put $\mathcal C=z_xp_xs_x=-z_yp_ys_y$ and
$\Delta_x=P/q-1$, $\Delta_y=S/q-1$. At a nonfresh state,
\[
 r=\mathcal C(\eta_x-\eta_y),\qquad
 h=\mathcal C(\eta_x\Delta_x-\eta_y\Delta_y),
\]
so
\[
 \ell:=h/r=\Delta_x+\frac{\Delta_x-\Delta_y}{J-1},
 \qquad 0<\ell<4P/q<10^5.
\]
Also $\gamma<10^5$ and $\zeta=1-S/q\in[1/2,1]$. In particular
$\sqrt H=\ell\sqrt Q$ without a sign ambiguity. Direct substitution gives
\begin{equation}\label{eq:bsur-qalpha}
 \frac Q\alpha=
 \frac{\kappa P}{1+\kappa\rho}(1-J^{-1})^2\frac{x^2}{y}.
\end{equation}
Its coefficient lies between $10^{-6}$ and $10^9$: for the lower bound use
$1-J^{-1}>1/4$ and
$(10^{-4})(1.02)/(2\cdot16)>10^{-6}$.
Define $W=x^2/(\omega y)$ and $\mathcal Z=Q/(\alpha g)$. We obtain
\begin{equation}\label{eq:bsur-gate-ratios}
 10^{-3}W\le\mathcal Z\le10^{14}W,\qquad
 1-G/g=(2\ell+\zeta)\mathcal Z\le10^6\mathcal Z,
 \qquad
 \frac{B_{\rm sur}}g=(\sqrt U+\gamma\sqrt{\mathcal Z})^2.
\end{equation}
The last expression retains the nonzero fresh contribution $U$; it is not the
corresponding expression for $K$.
An accepted nonfresh origin has $\mathcal Z\le1$ and $W\le10^3$. By
\eqref{eq:bsur-ratio-region}, its curvature ratio is below
$(1+10^5)^2<4\cdot10^{10}$. Thus its actual selected step satisfies
\begin{equation}\label{eq:bsur-step-floor}
 t\ge\min\{1/2,(32\cdot10^{10})^{-1}\}>10^{-12}.
\end{equation}
The fresh step obeys the same lower bound by~\eqref{eq:bsur-fresh-step}.

\paragraph{Invariance and finite exit.}
The initial ratios in~\eqref{eq:bsur-entry} are far below $10^{-20}$.
Its first actual Bsur step preserves the region by
\eqref{eq:bsur-multipliers}. Subsequently $y/\omega$ decreases and remains at
most $e^{-b}$. At any accepted nonfresh origin already in the region,
\[
 x/\omega=\sqrt{W(y/\omega)}\le\sqrt{10^3}e^{-b/2}.
\]
Its next update gives $x^+/\omega^+<10^4\sqrt{10^3}e^{-b/2}<10^{-20}$;
the other ratio and $\omega$ decrease. This proves the region by induction,
including the first failed endpoint. It does not require that endpoint to pass
the gate.

For each accepted step put
\[
 a_k=\log(F_P/F_\delta)>0,\qquad
 \beta_k=\log(F_\delta/F_S)>0.
\]
Their exact increments satisfy
\begin{equation}\label{eq:bsur-log-motion}
 \log W_{k+1}-\log W_k=2a_k+\beta_k,
 \qquad \beta_k\ge\log(1+.98t_k)\ge t_k/2.
\end{equation}
The latter follows from $F_\delta-F_S>.98t_k$ and $F_S\le1$.
Every accepted step therefore increases $\log W$ by more than
$5\cdot10^{-13}$, whereas every accepted origin has $W\le10^3$, including
$W_0=e^{-3b-2\varepsilon}$ at the fresh state. The epoch must end after finitely
many accepted steps $M$; a conservative explicit bound is
\[
 M\le2+\left\lceil2\cdot10^{12}
                  (3b+200+\log10^3)\right\rceil.
\]
At the first failed endpoint, $g>0$ and $G/g<1/2$, so
\[
 1/2<10^6\mathcal Z\le10^{20}W_{\rm end}.
\]
The preceding accepted origin has $W\le10^3$, and its final multiplier for W
is below $10^{10}$. Consequently
\begin{equation}\label{eq:bsur-endpoint}
 10^{-22}<W_{\rm end}\le10^{13},\qquad
 |\log W_{\rm end}|\le100.
\end{equation}
No current fit, approximate refresh threshold, or assumed exit time was used.

\subsection{Summable error over the actual variable-step prefix}
\label{sec:bsur-prefix}

Use the reference step $t_0=1/(4\delta)=25/99$ and define
\[
 a_\sigma(t)=\log\frac{1+t(\mathcal A_\sigma-1)}{1+t(\delta-1)},
 \qquad
 \beta(t)=\log\frac{1+t(\delta-1)}{1-t}.
\]
These are reference increments; no actual step is assumed to equal $t_0$.
Write $a_\sigma=a_\sigma(t_0)$ and $\beta=\beta(t_0)$.

\paragraph{A summable bound on $U-\delta$.}
Since $qg=\sum_iw_i(\lambda_i-q)^2$,
\[
 U-\delta/q=
 \frac{\sum_iw_i(\lambda_i-\delta)(\lambda_i-q)^2}{q^2g}.
\]
The auxiliary numerator is zero. Using $q^2g\ge.999^2\,10^{-5}\omega$,
$1-q\le.011\omega$, and the score bounds yields
\begin{equation}\label{eq:bsur-u-error}
 |U-\delta|\le\omega+2\cdot10^{17}x/\omega+2\cdot10^5y/\omega.
\end{equation}
For clarity, the three coefficients before rounding upward are bounded by
\[
 \frac{\delta(.011)}{.999}
 +\frac{.01(.011)^2}{.999^2\,10^{-5}}<1,\qquad
 \frac{11000^3}{.999^2\,10^{-5}}<2\cdot10^{17},\qquad
 \frac1{.999^2\,10^{-5}}<2\cdot10^5.
\]

\paragraph{The prefix and its actual steps.}
Let $W_{\rm cut}=10^{-40}$ and let $n$ be the first index with
$W_n>W_{\rm cut}$. It exists before or at the failed endpoint: initially W is
below the cutoff, the failed endpoint is above it, and all states below the
cutoff will pass the gate. At any prefix origin $k<n$,
\[
 \mathcal Z\le10^{-26},\qquad
 \eta:=1-G/g\le10^{-20},\qquad
 \gamma\sqrt{\mathcal Z}\le10^{-8}.
\]
At the fresh origin these upper bounds hold with $\mathcal Z=\eta=0$, without
defining $h/r$ there. Put
\[
 V=B_{\rm sur}/g=(\sqrt U+\gamma\sqrt{\mathcal Z})^2.
\]
Weakening the regional upper bound to $U<3/2$ gives
$3/4<V<(5/4+10^{-8})^2<7/4$. The prefix gate passes, and the half-step cap is
inactive. Thus the actual step obeys
\begin{equation}\label{eq:bsur-prefix-step}
 t_k=\frac{1-\eta}{4V},\qquad 1/8<t_k<1/3.
\end{equation}
Its discrepancy from the reference step is bounded by
\begin{align*}
 |t_k-t_0|
 &\le\frac{\eta}{4V}+\frac{|V-\delta|}{4\delta V}\\
 &\le\frac\eta3+\frac12|U-\delta|
       +\gamma\sqrt U\sqrt{\mathcal Z}
       +\frac12\gamma^2\mathcal Z.
\end{align*}
The cross term is at most $(5/4)10^{12}\sqrt W$. The first and last terms
together are at most $(10^6/3+10^{10}/2)10^{14}W$.
Since $W\le10^{-40}$ implies $W\le10^{-20}\sqrt W$, these estimates give
\begin{equation}\label{eq:bsur-step-error}
 |t_k-t_0|\le |U-\delta|+2\cdot10^{12}\sqrt{W_k}.
\end{equation}

\paragraph{Uniform sums, without multiplying by the epoch length.}
By~\eqref{eq:bsur-log-motion} and~\eqref{eq:bsur-prefix-step}, each prefix
step raises $\log W$ by at least $1/16$. Summing backward from
$W_{n-1}\le10^{-40}$, and using the forward contraction of $y/\omega$, gives
\begin{equation}\label{eq:bsur-prefix-sums}
 \begin{aligned}
 \sum_{k<n}\sqrt{W_k}&\le33\cdot10^{-20},&
 \sum_{k<n}\omega_k&\le(8000/9)\omega_s,\\
 \sum_{k<n}y_k/\omega_k&\le17e^{-b},&
 \sum_{k<n}x_k/\omega_k&\le33\cdot10^{-20}e^{-b/2}.
 \end{aligned}
\end{equation}
Indeed $1/(1-e^{-1/32})<33$ bounds the first geometric sum. The second uses
$\omega_{k+1}\le(1-9/8000)\omega_k$, and the third uses
$1/(1-e^{-1/16})<17$. These denominator bounds follow from
$e^{-u}\le1/(1+u)$. The fourth estimate uses the identity
$x/\omega=\sqrt{W(y/\omega)}$ and $y/\omega\le e^{-b}$.

Combining~\eqref{eq:bsur-u-error}, \eqref{eq:bsur-step-error}, and
\eqref{eq:bsur-prefix-sums}, for every $b\ge1000$,
\begin{equation}\label{eq:bsur-total-step-error}
 \begin{aligned}
 \sum_{k<n}|t_k-t_0|
 &\le(8000/9)10^{-6}
 +(2\cdot10^{17})(33\cdot10^{-20})e^{-500}\\
 &\quad +(2\cdot10^5)(17)e^{-1000}
 +(2\cdot10^{12})(33\cdot10^{-20})
 <10^{-3}.
 \end{aligned}
\end{equation}
The last inequality remains strict on replacing the exponentials by the larger
$2^{-500},2^{-1000}$. Also,
\begin{equation}\label{eq:bsur-total-q-error}
 \sum_{k<n}|q_k-1|\le.011\sum_{k<n}\omega_k<10^{-5}.
\end{equation}
These bounds control the entire prefix independently of its length.

\subsection{The logarithmic epoch map and the original finite start}
\label{sec:bsur-map}

The last prefix step has $W_n\le10^{10}W_{\rm cut}=10^{-30}<1$.
Its accumulated log-W motion is at least $n/16$, while
$-\log W_0=3b+2\varepsilon$. Therefore
\[
 n<16(3b+200)\le100b.
\]
The fixed-cache residual errors in~\eqref{eq:bsur-scores} satisfy
\[
 |P-\mathcal A_\sigma|\le2.2\cdot10^8e^{-b+100},\qquad
 S\le1.1\cdot10^8e^{-2b+200}.
\]
Consequently
\begin{equation}\label{eq:bsur-total-residual-error}
 n|P-\mathcal A_\sigma|
 \le100b(2.2\cdot10^8)e^{-b+100}<1,
 \qquad nS\le100b(1.1\cdot10^8)e^{-2b+200}<1.
\end{equation}
Both upper bounds decrease for $b\ge1000$; at $b=1000$ even the larger powers
$2^{-900},2^{-1800}$ suffice. Multiplication by the prefix length is legitimate
here because that length has been bounded and the cache errors decay
exponentially in the entry depth.

The exact variable-step logarithmic increments are
\[
 \begin{aligned}
 a_k&=\log[q_k(1-t_k)+t_kP]-\log[q_k(1-t_k)+t_k\delta],\\
 \beta_k&=\log[q_k(1-t_k)+t_k\delta]-\log[q_k(1-t_k)+t_kS].
 \end{aligned}
\]
For $q\in[.999,1]$, $t\in[1/8,1/3]$, and the intervening nonnegative score
ranges, either log ratio has derivatives bounded in absolute value by 16 in t,
3 in q, and 1 in its informative score P or S. For example,
\[
 \left|\partial_t\log(q(1-t)+tU)\right|
 \le\max\{1/t,1/(1-t)\}\le8,\qquad
 0\le\partial_q\log(q(1-t)+tU)\le1/q.
\]
The score derivative is at most $t/[q(1-t)]<1$.
The same bounds hold on the comparison paths to $q=1$, $t=t_0$,
$P=\mathcal A_\sigma$, and $S=0$.
Equations~\eqref{eq:bsur-total-step-error}--\eqref{eq:bsur-total-residual-error}
therefore give
\begin{equation}\label{eq:bsur-prefix-log-error}
 \left|\sum_{k<n}a_k-na_\sigma\right|<2,\qquad
 \left|\sum_{k<n}\beta_k-n\beta\right|<2.
\end{equation}
Thus the actual varying prefix has been controlled, rather than replaced by a
fixed-step recurrence.

\paragraph{The unrestricted tail and endpoint reset.}
No step approximation is used after the cutoff. Positivity of both logarithmic
increments and~\eqref{eq:bsur-endpoint} imply
\[
 \sum_{k=n}^{M-1}(2a_k+\beta_k)
 =\log(W_{\rm end}/W_n)<53\log10<159.
\]
The tail contributes less than $79.5$ to the a sum and $159$ to the beta sum,
regardless of the number and sizes of its actual steps. With
$X=-\log(x/\omega)$ and $Y=-\log(y/\omega)$, we consequently have
\[
 X_{\rm end}=2b+\varepsilon-na_\sigma+E_X,\qquad
 Y_{\rm end}=b+n\beta+E_Y,\qquad |E_X|,|E_Y|<200.
\]
Set $b'=X_{\rm end}$ and
$\varepsilon_{\rm next}=Y_{\rm end}-2X_{\rm end}=\log W_{\rm end}$.
Eliminating n gives
\begin{equation}\label{eq:bsur-exact-map}
 \begin{aligned}
 b'={}&C_\sigma b
 +\frac{\beta}{2a_\sigma+\beta}\varepsilon
 -\frac{a_\sigma}{2a_\sigma+\beta}\varepsilon_{\rm next}\\
 &+\frac{a_\sigma}{2a_\sigma+\beta}E_Y
 +\frac{\beta}{2a_\sigma+\beta}E_X,
 \qquad C_\sigma=\frac{a_\sigma+2\beta}{2a_\sigma+\beta}.
 \end{aligned}
\end{equation}
Every error coefficient has modulus below one. Hence the actual Bsur epoch obeys
\begin{equation}\label{eq:bsur-map-bounds}
 |b'-C_\sigma b|<1000,\qquad |\varepsilon_{\rm next}|\le100.
\end{equation}
The first bound follows already with 600. The second is a new endpoint bound,
not an accumulation of the old epsilon parameters. If $b'\ge1000$, then
$2b'+\varepsilon_{\rm next}>b'$, so the informative rows really exchange
their small and large roles when the next epoch is relabeled.

\paragraph{A finite two-epoch margin.}
At $t_0=25/99$ the constants simplify to
\[
 a_{\mathrm H}=\log(1030597/395),\qquad
 a_{\mathrm L}=\log(3990301/3950000),\qquad
 \beta=\log(395/296).
\]
The following strict finite bounds suffice:
\[
 a_{\mathrm H}<7.87,\qquad a_{\mathrm L}<.01016,\qquad \beta>.2885.
\]
For the first, the positive Taylor polynomial
$\sum_{j=0}^{20}(7.87)^j/j!$ exceeds $1030597/395$ and is below $e^{7.87}$.
For the second, use $\log(1+x)\le x-x^2/2+x^3/3$ at
$x=40301/3950000$. For the third, put $z=99/691$ and use
\[
 \beta=2\operatorname{arctanh}z
 >2(z+z^3/3)=95188104/329939371>.2885.
\]
The function $C(a,\beta)=(a+2\beta)/(2a+\beta)$ decreases in a and increases
in beta. Thus
\begin{equation}\label{eq:bsur-map-expansion}
 C_{\mathrm H}C_{\mathrm L}>
 \frac{7.87+.577}{15.74+.2885}
 \frac{.01016+.577}{.02032+.2885}
 =\frac{165324684}{164997379}>1.0019,
 \qquad 1/2<C_\sigma<2.
\end{equation}
This is a smaller expansion than for K, and requires its own finite entry check.

Suppose an H-type epoch starts with $b_H\ge5\cdot10^6$. Its next parameter
satisfies $b_L>b_H/2-1000\ge2499000>1000$, so the actual role exchange is
valid and the next L-type epoch satisfies~\eqref{eq:bsur-entry}.
After that L epoch, the lower bound $C_{\mathrm L}>1/2$ first gives a new
parameter greater than $1248500$, independently verifying the second role
exchange. Combining the two maps more sharply gives
\begin{equation}\label{eq:bsur-global-induction}
 b_{H,\rm next}>1.0019b_H-3000\ge1.001b_H.
\end{equation}
The final inequality has at least 1500 of slack at $b_H=5\cdot10^6$.
The parameter error in b has therefore been charged at both epochs; the
endpoint epsilon bound is not used to suppress it.

For the unchanged individual-row uniform initialization,
\[
 x_1/\omega=N^{-2},\qquad x_2/\omega=N^{-1},\qquad
 b_0=10^7\log2>5\cdot10^6,\qquad
 \varepsilon_0=0,\qquad \omega_0<1/N<10^{-6}.
\]
It enters the H-type induction exactly. All subsequent H parameters increase
geometrically and every intervening L parameter remains above 1000.
Equations~\eqref{eq:bsur-entry}--\eqref{eq:bsur-map-bounds} can therefore be
iterated indefinitely, with the original finite N and no modified initialization.
Each epoch is finite, so this yields infinitely many actual exact refreshes.

\subsection{The value limit and the scope of the conclusion}
\label{sec:bsur-conclusion}

Every finite multiplier is at least $1-t\ge1/2$, so all individual weights stay
positive and every cache is legal. There is no fresh zero-g termination.
The proved surrogate safeguard gives
\[
 v(w+td)-v(w)\ge tG/2\ge tg/4>0
\]
at every accepted update. The auxiliary mass decreases at every update and
contracts by a factor below .99775 at each of the infinitely many fresh steps.
Thus $\omega\to0$; the invariant ratios imply $x_1,x_2\to0$ and $w_0\to1$.
By~\eqref{eq:bsur-value} and evaluation at $\theta=0$,
\[
 w_0\le v(w)\le w_0+\delta\omega+B(x_1+x_2)
 =1-.01\omega+.01(x_1+x_2)<1.
\]
This proves~\eqref{eq:bsur-claim}.

The construction is analytic; no enormous matrix or trajectory is claimed to
have been executed. It gives a second specified uncapped rule with the same
suboptimal limit, not an identity of its iterates with those of the K rule.
The conclusion concerns the unchanged G margin and this Bsur step. It does not
assert failure of every auxiliary-function method, the capped Bsur policy,
corrected-fit methods, bounded-delay policies, or uniform starts on distinct
rows. Nor does changing the coefficient make the two trajectories have the same
total work: forming $R_2$ requires additional scalar work, and their update and
refresh counts need not agree. Finally $A_w\to0$, while its scalar spectral
condition number remains one; the full information matrix becomes singular.

\section{The uncapped equal-magnitude nuisance class}
Suppose $p=1$ and $z_i\in\{-c,0,c\}$ for some $c>0$. The original uncapped relative-margin rule converges from every positive initial design. This special statement does not extend to the unequal-magnitude counterexample above.

The finite-refresh branch was proved earlier. With infinitely many refreshes, the epoch argument gives $q_k\to\bar v>0$, and the lower bound on fresh steps gives $\sum_k t_k=\infty$, with $0<t_k\le1/2$.
For real $b,u$, $q>0$ and $0<t\le1/2$,
\begin{equation}\label{eq:pair-product}
 [1-t+t(b+u)^2/q][1-t+t(b-u)^2/q]\ge1+2t(b^2/q-1).
\end{equation}
Indeed, put $\beta=b^2/q$, $x=u^2/q$. The product is
$F(x)=(1-t+t\beta+tx)^2-4t^2\beta x$.
If its unconstrained minimizer $x_*=\beta-(1-t)/t$ is nonpositive, its minimum on $x\ge0$ is $(1+t(\beta-1))^2$, proving the inequality.
Otherwise its minimum is $4t\beta(1-t)$; subtracting the right side gives
$(1-2t)(2t\beta-1)\ge0$, since $\beta>(1-t)/t$.

Orient nonzero-nuisance rows by $\operatorname{sign}(z_i)$ and put
\[
 b_i=\operatorname{sign}(z_i)a_i,\quad
 D_*=\tfrac12(\max_{z_i\ne0}b_i-\min_{z_i\ne0}b_i),\quad
 A_0=\max_{z_i=0}a_i^2,
\]
using $A_0=0$ if that set is empty. Then $v_*=\max\{A_0,D_*^2\}$:
choose the coefficient minus the midpoint of the extreme $b_i$, divided by $c$, for an upper bound; a zero-nuisance singleton attains $A_0$ when present. If $D_*>0$, the distinct extreme pair with equal masses attains $D_*^2$; if $D_*=0$, its lower bound is trivial, and no pair (or self-pair) is required.
Only in the case $\bar v<D_*^2$ is a pair argument needed; then $D_*>0$. After negating one residual in that pair, its residuals are $D_*+u_k,D_*-u_k$ for every cache, even if fits oscillate.
If $\bar v<D_*^2$, eventually $D_*^2/q_k\ge1+\eta$ for some $\eta>0$.
Equation~\eqref{eq:pair-product} multiplies their positive weight product by at least $1+2\eta t_k$. Its logarithm would diverge, since
$\log(1+2\eta t_k)\ge2\eta t_k/(1+\eta)$, contradicting the product bound $1/4$.
A zero-nuisance row similarly excludes $\bar v<a_i^2$ through growth of its weight.
Hence $\bar v\ge\max\{A_0,D_*^2\}=v_*$, proving equality without a weight floor or convergence of cached fits.

\section{A sharp high-accuracy penalty of the refresh-count cap}
Consider the capped $K$ rule on $(1,-1),(1,1)$ with first current mass $x=(1-y)/2$ and cache mass $s_1=(1-a)/2$, where $0<y\le a<1$. The cached residuals are $(1+a,1-a)$, $A_s=1$, and substitution gives
\[
 q=1+a^2-2ay,\quad d_1=-d_2=d=\frac{a(1-y^2)}q,\quad g=4ad,
\]
\[
 r=y-a,\quad h=-2d,\quad\alpha=\frac{1+y}{1+a},\quad
 \gamma=\frac{2d}{1-y},\quad\zeta=\frac{2d}{1+y},
\]
\[
 \frac Gg=1-\frac{1+a}{1+y}
 \left(\frac{a-y}{a}+\frac{(a-y)^2}{2a(1+y)}\right),\quad
 K=\frac{4a^2(1+y)(1+a)(1+a-2y)^2}{q^2}.
\]
For $z=y/a$, call the first ratio $F_a(z)$. The gate satisfies
\[
 2(1+az)^2(F_a(z)-1/2)
 =(a+2a^2)z^2+(2+4a)z-(1+3a+a^2)=:P_a(z).
\]
On $[0,1]$, $P_a'>0$, $P_a(1/2)=-3a/4-a^2/2<0$ and $P_a(1)=(1+a)^2>0$.
Its unique zero $z_*(a)\in(1/2,1)$ is the exact retained-cache threshold, with $z_*(a)=1/2+3a/8+O(a^2)$.

The true directional derivative is $4yd$, so $G\le4yd$.
Also $K=(2d+\gamma(a-y))^2/\alpha\ge4d^2$. Therefore
\[
 t\le G/(4K)\le y/(4d),\qquad y/2\le y^+=y-2td<y.
\]
No finite step hits or crosses the optimum. Capped convergence gives $y\to0$.
There must be infinitely many refreshes: a fixed final positive $a$ would contradict $y/a\ge z_*(a)>1/2$.
Let $a_j$ be the cache bias after refresh $j$, including the initial $j=0$, and $n_j$ its number of accepted steps.
Uniformly for $z\in[0,1]$, $g/a^2\to4$ and $K/a^2\to4$ as $a\to0$.
On the accepted set $G/g\ge1/2$, the uncapped safe step therefore has a fixed positive lower bound in all sufficiently late epochs. The cap $1/(j+1)$ eventually determines every step in those epochs.
Their exact recurrence is
\[
 z^+=z-\frac2{j+1}
       \frac{1-a_j^2z^2}{1+a_j^2-2a_j^2z}.
\]
The last factor converges uniformly to one and the threshold tends to $1/2$.
Exit overshoots it by at most $O(1/j)$, giving
\[
 z_{\rm exit}\to1/2,\qquad n_j/(j+1)\to1/4,\qquad
 a_{j+1}/a_j=z_{\rm exit}\to1/2.
\]
At the start of epoch $J$, $N_J=\sum_{j<J}n_j$ satisfies
$N_J/J^2\to1/8$ and $\log a_J/J\to-\log2$.
Within that epoch $N=N_J+O(J)$ and $\log y=\log a_J+O(1)$, because $y/a_J$ stays bounded away from zero.
The efficiency gap is $e_N=y_N^2$, so
\[
 \frac{\log e_N}{\sqrt N}\longrightarrow-4\sqrt2\log2.
\]
Strictly decreasing gaps permit inversion:
\[
 N_\epsilon\sim\frac{\log^2(1/\epsilon)}{32\log^2 2},\qquad
 J_\epsilon\sim\frac{\log(1/\epsilon)}{2\log2}.
\]
For the uncapped rule every fresh state $y=a$ has $d=a$, $g=G=K=4a^2$, so $t=1/4$ and $y^+=a/2$.
The strict inequality $P_a(1/2)<0$ forces a refresh before the next step.
Thus $y_N=2^{-N}y_0$ exactly.
Fresh half-affine averages the swapped masses, while half-power gives equal unnormalized masses $2x(1-x)$; both are optimal in one update.
Repeating the $+1$ row gives a full-rank uniform three-row initialization with aggregate $x_0=1/3$.
This is a duplicate-row example, not a distinct-site uniform-start result.

\paragraph{Generic-operation translation.}
For a completed nonstationary prefix of $N\ge1$ accepted steps and $J$ actual refreshes, stop counting at the final accepted update.
There are $N+J$ gates, of which $J+1$ are fresh and $N-1$ stale.
The recorded generic implementation consequently has
\[
 F=J+1,\quad H=2N+J-1,\quad M=2N+3J+1,\quad T=2(N+J).
\]
Here $F$ counts each assembly, factorization and cache contrast solve; $H$ old-factor solve pairs; $M$ feature passes; $T$ scalar reduction passes.
Fresh gates pay one old solve and stale gates two, including rejected gates; each cache also pays two feature passes.
Let fixed nonnegative prices be $B_0,h,\mu,\nu$, with $B_0$ the sum of assembly/factor/contrast prices. Then $C_K=B_0F+hH+\mu M+\nu T$.
Writing $L_\epsilon=\log(1/\epsilon)$, for $h+\mu+\nu>0$,
\[
 C_{\rm cap}\sim\frac{h+\mu+\nu}{16\log^2 2}L_\epsilon^2,\qquad
 C_{\rm uncapped}\sim\frac{B_0+3h+5\mu+4\nu}{2\log2}L_\epsilon.
\]
Their cost ratio diverges, while the leading refresh counts agree.
These are operation prices, not measured hardware coefficients.
Special scalar formulas may remove solves for both implementations, but not the accepted-update count separation.
Additional priced nuisance inner products are considered separately below; no hardware speedup follows.

\section{First-hit cost regions including preconditioned correction}
This is a static calculation from saved runs, with no new trajectory, timing or LP execution. Every row uses that method's first recorded offline efficiency at least $0.999999$, not a later matched-budget iterate. The target is not an online decision input.
Extend the preceding price model to
\[
 C=B_0F+hH+\mu M+\nu T+\xi D,
\]
where $D$ counts explicit nuisance-vector inner products outside solves and feature matvecs. These length-$p$ dots differ from length-$m$ node reductions. This symbolic model still does not measure all physical FLOPs, memory traffic or hardware overhead.
Old $K$ and $B_{\rm sur}$ gates have $D=H$. A correction computes two dots for its line minimizer and one more only when its corrected residual is nonzero. If $N_c$ corrections and $N_y$ nonzero corrected residuals occur, $H=N_c+N_y$ and $D=2N_c+N_y=H+N_c$.
The saved primitive counts record no zero-initial-residual stale skips on the three runs, so the two attained correction first hits have $N_c=N-1$. Fresh half-power/half-affine have $D=0$. Dots internal to old-factor triangular solves are already paid in $H$; those internal to a fresh contrast solve are paid in $B_0$.

\begin{center}\small
\begin{tabular}{llrrrrrr}
\toprule
Case & Method & Steps & $F$ & $H$ & $M$ & $T$ & $D$\\
\midrule
$(1,2)$ & Half-power &1&1&0&2&1&0\\
 & Half-affine &2&2&0&4&2&0\\
 & Uncapped $K$ &8&8&22&38&30&22\\
 & Capped $K$ &9&6&22&34&28&22\\
 & Capped $B_{\rm sur}$ &10&6&24&36&45&24\\
 & Correction &9&1&13&26&27&21\\
\midrule
Cubic & Half-power &203&203&0&406&203&0\\
 & Half-affine &206&206&0&412&206&0\\
 & Uncapped $K$ &210&106&524&736&630&524\\
 & Correction &879&12&1756&2658&2670&2634\\
\bottomrule
\end{tabular}
\end{center}
The old-fit capped cubic methods miss the target, and all six methods miss it on the difficult boundary case; those rows are excluded from first-hit comparisons. The $H$ column is old solves, not the total-solve column in the main experimental table.

Half-power is the unique Pareto-efficient method on $(1,2)$; correction itself dominates all three old-fit safeguards there. For the cubic case, half-power dominates half-affine, leaving half-power $P$, uncapped $U$ and correction $C$. Define
\[
 T_{UP}=\frac{524h+330\mu+427\nu+524\xi}{97},\qquad
 T_{CU}=\frac{1232h+1922\mu+2040\nu+2110\xi}{94}.
\]
Every coefficient of $T_{CU}$ exceeds its counterpart in $T_{UP}$. For strictly positive prices $T_{UP}<T_{CU}$, and the correction/half-power threshold is
\[
 T_{CP}=\frac{97T_{UP}+94T_{CU}}{191}.
\]
Subtracting the table's resource vectors proves that $P$ is cheapest for $B_0<T_{UP}$, $U$ for $T_{UP}<B_0<T_{CU}$, and $C$ for $B_0>T_{CU}$, with ties at the boundaries. For example, correction beats half-power precisely when
$191B_0>1756h+2252\mu+2467\nu+2634\xi$.
Setting $\xi=0$ recovers the four-price projection; it deliberately omits a separately priced resource rather than claiming exhaustive arithmetic accounting.

For generic $K$ prefixes the extension replaces $h$ by $h+\xi$ in the preceding cost asymptotics, since $D=H$. Thus its logarithmic-squared separation persists for fixed nonnegative prices with $h+\xi+\mu+\nu>0$. It does not follow if only cache construction is priced, or if prices change with requested accuracy.
All first-hit determinations use separately charged offline truth. No certified online stop or physical price calibration is added by this calculation.

\paragraph{A direct inexact-gradient comparison.}
On this cap-cost instance the late epoch lengths satisfy $n_j/(j+1)\to1/4$, and each update of that epoch has $t_k=1/(j+1)$. Therefore
\[
 \sum_{k\text{ in epoch }j}t_k^2
 =\frac{n_j}{(j+1)^2}\sim\frac1{4(j+1)}.
\]
The sum over all accepted steps diverges (eventually each epoch contributes at least $1/[8(j+1)]$). Thus identifying $\gamma_k=t_k$ does not meet the square-summability hypothesis of Bertsekas and Tsitsiklis~\cite[Proposition~1]{bertsekas2000}. That proposition additionally assumes a globally Lipschitz gradient on an unconstrained domain and uniform direction/error bounds; these are not supplied by the simplex slope gate alone. The comparison does not rule out other transformations or convergence theorems.

\section{Specific interfaces with general surrogate convergence theorems}
\label{app:mm-interface}

This section makes the prior-method comparison precise; it adds no new update
rule or contribution claim. For a positive exact anchor $s$, put
\[
 b_i(s)=s_ip_i(s)/v_s,\qquad
 L_s(w)=\left(\sum_i b_i(s)^2/w_i\right)^{-1}\le v(w),
 \qquad L_s(s)=v_s.
\]
The corresponding minimization majorant for $f=-v$ is $u(w,s)=-L_s(w)$.
The SUM algorithm of Razaviyayn, Hong and Luo~\cite[Assumption~1 and
Theorem~1, arXiv v1, pp.~4--8]{razaviyayn2013} exactly minimizes the
current model. Its hypotheses include value contact, majorization,
directional first-order agreement (A3), and joint continuity (A4).
It permits nonsmooth objectives. Mairal's first-order-surrogate
framework~\cite[Definition~2.2, Algorithm~1 and Proposition~2.5]{mairal2015}
also exactly minimizes the chosen model and requires a common finite
Lipschitz constant for the gradients of the surrogate errors.
We check these specific interfaces rather than excluding nonsmooth
optimization theories as a class.

\paragraph{The fresh affine step is not the exact covariance-model optimizer.}
Cauchy--Schwarz gives the full-simplex maximizer of $L_s$ as
$w_i^*=s_i|p_i|/\sum_j s_j|p_j|$ when every $p_i\ne0$.
The fresh affine step instead has multiplier
$w_i^+/s_i=1-t+tp_i^2/v_s$. In our counterexample the four positive squared
scores obey $S<\delta<1<P$. Equality of the two updates would require all
four different numbers $z_i=|p_i|$ to solve
\[
 (t/v_s)z_i^2-Cz_i+(1-t)=0
\]
for one $C>0$. A nonzero quadratic has at most two roots, and $t>0$.
Thus the specified fresh steps, capped or uncapped, are not exact optimizers
of this particular model. This does not negate the model-improvement
monotonicity argument of Yu~\cite{yu2010} or exclude other inexact-MM rules.

\paragraph{Directional agreement at the bad boundary.}
In aggregate coordinates write $e_0=(1,0,0,0)$ and let $u_*$ be the informative
pair optimum with masses $1/(R+1),R/(R+1)$. The value formula gives
\[
 v((1-h)e_0+hu_*)=1+(B-1)h,\qquad v'(e_0;u_*-e_0)=B-1=.01.
\]
The natural boundary estimator is $b(e_0)=e_0$, so
$L_{e_0}(w)=w_0$ and $L'_{e_0}(e_0;u_*-e_0)=-1$.
Thus this extension violates A3. In particular $-v$ is not directionally
stationary at $e_0$; referring only to single-row transfer directions would
miss the improving pair direction.

\paragraph{No jointly continuous, contacting extension of this model.}
For small $h>0$, take the two positive designs
\[
 \begin{aligned}
 s(h)&=(1-h-h^2,h^2,h/(R+1),Rh/(R+1)),\\
 w(h)&=(1-h-h^2-h^3,h^2,h^3,h).
 \end{aligned}
\]
Both tend to $e_0$. At the anchor the exact fit is zero and
$v(s)=1+(B-1)h+(\delta-1)h^2\to1$. The first informative coefficient
$b_1(s)=\sqrt B h/[(R+1)v(s)]$ yields
\[
 \sum_i\frac{b_i(s(h))^2}{w_i(h)}
 \ge\frac{B}{(R+1)^2v(s(h))^2h}\longrightarrow\infty,
 \qquad L_{s(h)}(w(h))\longrightarrow0.
\]
Contact would require $L_{e_0}(e_0)=1$. Hence any extension agreeing with
the interior formula and retaining contact fails joint continuity A4.
These are test paths for a model hypothesis, not purported algorithm iterates.

\paragraph{Error smoothness along actual fresh anchors.}
Let $s=(s_0,\omega,x_1,x_2)$ be an actual fresh anchor of the Bsur
counterexample. Write $\epsilon=1-s_0$, $r=(\omega,x_1,x_2)/\epsilon$,
and $w(\alpha)=(1-\alpha,\alpha r)$. Scaling every non-base mass equally
leaves the minimizing nuisance fit unchanged. If $c$ is the value of the
normalized non-base design, then
\[
 v(w(\alpha))=1+(c-1)\alpha,\qquad v_s=1+(c-1)\epsilon.
\]
The invariant $x_1+x_2\le2\cdot10^{-20}\omega$ gives
\[
 c\ge\frac{\delta\omega}{\epsilon}>.98,\qquad
 c\le\frac{\delta\omega+B(x_1+x_2)}{\epsilon}<1,
 \qquad 0<\epsilon<2\cdot10^{-6}.
\]
For the error $h_s=v-L_s$, the exact model along this local test ray is
\[
 L_s(w(\alpha))=\frac{v_s^2}{D(\alpha)},\qquad
 D(\alpha)=\frac{(1-\epsilon)^2}{1-\alpha}
                 +\frac{c\epsilon^2}{\alpha}.
\]
At $\alpha=\epsilon$, we have $D=v_s$, $D'=1-c$, and
$D''=2/(1-\epsilon)+2c/\epsilon$. Differentiating twice gives
\[
 \left.\frac{d^2}{d\alpha^2}h_s(w(\alpha))\right|_{\alpha=\epsilon}
 =\frac2{1-\epsilon}+\frac{2c}{\epsilon}
      -\frac{2(c-1)^2}{v_s}
 >\frac{1.96}{\epsilon}-.001>\frac1\epsilon.
\]
Since $\|w'(\alpha)\|^2=1+\|r\|^2\le2$, any local Lipschitz constant
for $\nabla h_s$ must exceed $1/(2\epsilon)$. This also holds before
aggregation: the two disjoint groups are scaled according to probability
vectors, whose squared norms sum to at most two. The actual fresh anchors
have $\epsilon\to0$, so there is no common finite error-smoothness constant
for this canonical model along that sequence. The ray is a local curvature
test at each actual anchor, not the update trajectory itself.

\paragraph{What the comparison does and does not exclude.}
A fixed old model generally lacks current contact, but the Bsur local proof
can use a latent exact current-contact model without computing its anchor.
The exact-optimizer and boundary conditions above remain separate obstacles;
absence of an executed current solve is not absence of an analysis model.
Mairal's incremental Algorithm~2 and Proposition~3.1 require randomly refreshed
component surrogates and exact minimization of their average. They are not
identified with our gate-triggered whole-cache refresh and affine partial
step; unbounded component age alone would not distinguish them.
Conversely, on a fixed positive-weight-floor simplex, exact minimization of
the current covariance model satisfies the SUM interface and yields a
restricted-domain optimum because $-v$ is convex. These concrete comparisons
neither rule out every surrogate representation or convergence theorem nor
constitute an absolute priority claim.


\section{Reproduction, limitations, and provenance}
The complete internal package contains the main and supplementary \LaTeX{} sources, bibliography, a build script, exact and numerical verification code, machine-readable results, an explicit claim--evidence table, a literature comparison, a failure log, and SHA256 manifests.
Run the documented experiment entry point to regenerate the diagnostics, and the build script to regenerate both PDFs. Exact rational identities establish only the displayed examples; seeded floating-point stress checks detect implementation inconsistencies but do not prove the theorem or validate roundoff safety.

The scalar statistical target, exact least-squares projection identity, concavity, fresh monotonicity, and feasible-dual reasoning are established background. The unpublished R11 source supplied the delayed-design and tangent-gap starting point. The present proof is self-contained and imports none of the protected spectral-memory, joint-root, or information-budget results from other manuscripts. Detailed internal file-level attribution is kept outside the anonymous manuscript in the source audit.

This draft establishes cached-residual safe steps, global capped convergence and actual failure of the uncapped relative-margin rule. Useful total-work gains and floating-point reliability are not established by these exact-arithmetic results. The general upper bound is too conservative to settle efficiency; the exact cap-cost obstruction and stronger comparators show why factor-count savings alone are insufficient. The model delays computations of deterministic design sensitivities, not experimental outcomes; no new statistical concentration theorem is claimed.

\clearpage
\begin{thebibliography}{99}
\bibitem{yu2010} Yaming Yu. Monotonic convergence of a general algorithm for computing optimal designs. \emph{The Annals of Statistics}, 38(3):1593--1606, 2010. \href{https://doi.org/10.1214/09-AOS761}{doi:10.1214/09-AOS761}. Preprint \href{https://arxiv.org/abs/0905.2646v4}{arXiv:0905.2646v4}.
\bibitem{zhao2025} Renbo Zhao. An Optimization Perspective on the Monotonicity of the Multiplicative Algorithm for Optimal Experimental Design. Preprint version used here: \href{https://arxiv.org/abs/2508.07074v1}{arXiv:2508.07074v1}, 2025. Publisher-assigned issue: \emph{Operations Research Letters}, 69, article 107511, November 2026. \href{https://doi.org/10.1016/j.orl.2026.107511}{doi:10.1016/j.orl.2026.107511}.
\bibitem{sagnol2023} Guillaume Sagnol and Luc Pronzato. Fast Screening Rules for Optimal Design via Quadratic Lasso Reformulation. \emph{Journal of Machine Learning Research}, 24(307):1--32, 2023. \url{https://jmlr.org/papers/v24/22-1305.html}.
\bibitem{woodruff2024} David P. Woodruff and Taisuke Yasuda. John Ellipsoids via Lazy Updates. \emph{Advances in Neural Information Processing Systems}, 2024. Archival preprint \href{https://arxiv.org/abs/2501.01801v1}{arXiv:2501.01801v1}, 2025.
\bibitem{macedo2026} Humberto Gimenes Macedo and Lu\'is Felipe Bueno. An inexact first-order descent method with general directions: theory and applications to DE-constrained optimization. \emph{Computational and Applied Mathematics}, 45, article 413, 2026. \href{https://doi.org/10.1007/s40314-026-03794-2}{doi:10.1007/s40314-026-03794-2}.
\bibitem{harman2008} Radoslav Harman and Tom\'a\v{s} Jur\'ik. Computing c-optimal experimental designs using the simplex method of linear programming. \emph{Computational Statistics \& Data Analysis}, 53(2):247--254, 2008. \href{https://doi.org/10.1016/j.csda.2008.06.023}{doi:10.1016/j.csda.2008.06.023}.
\bibitem{bental1990} Aharon Ben-Tal and Marc Teboulle. A geometric property of the least squares solution of linear equations. \emph{Linear Algebra and its Applications}, 139:165--170, 1990. \href{https://doi.org/10.1016/0024-3795(90)90395-S}{doi:10.1016/0024-3795(90)90395-S}.
\bibitem{derezinski2018} Micha\l{} Derezi\'nski and Manfred K. Warmuth. Reverse Iterative Volume Sampling for Linear Regression. \emph{Journal of Machine Learning Research}, 19(23):1--39, 2018. \url{https://jmlr.org/papers/v19/17-781.html}.
\bibitem{lupong2013} Zhaosong Lu and Ting Kei Pong. Computing Optimal Experimental Designs via Interior Point Method. \emph{SIAM Journal on Matrix Analysis and Applications}, 34(4):1556--1580, 2013. \href{https://doi.org/10.1137/120895093}{doi:10.1137/120895093}. Preprint \href{https://arxiv.org/abs/1009.1909v3}{arXiv:1009.1909v3}.
\bibitem{freund1997} Yoav Freund and Robert E. Schapire. A Decision-Theoretic Generalization of On-Line Learning and an Application to Boosting. \emph{Journal of Computer and System Sciences}, 55(1):119--139, 1997. \href{https://doi.org/10.1006/jcss.1997.1504}{doi:10.1006/jcss.1997.1504}. Author manuscript: \url{https://www.schapire.net/papers/FreundSc95.pdf}.
\bibitem{richtarik2008} Peter Richt\'arik. Simultaneously solving seven optimization problems in relative scale. Technical report, December 2008; posted January 2009. \url{https://optimization-online.org/2009/01/2185/}.
\bibitem{cohen2019} Michael B. Cohen, Ben Cousins, Yin Tat Lee, and Xin Yang. A near-optimal algorithm for approximating the John Ellipsoid. \emph{Proceedings of the Thirty-Second Conference on Learning Theory}, PMLR 99:849--873, 2019. \url{https://proceedings.mlr.press/v99/cohen19a.html}. Revised preprint \href{https://arxiv.org/abs/1905.11580v2}{arXiv:1905.11580v2}, 2020.
\bibitem{torsney1981} Bernard Torsney. \emph{Algorithms for a constrained optimization problem with applications in statistics and optimum design}. PhD thesis, University of Glasgow, 1981. \url{https://theses.gla.ac.uk/1088/}.
\bibitem{pukelsheim1991} Friedrich Pukelsheim and Ben Torsney. Optimal Weights for Experimental Designs on Linearly Independent Support Points. \emph{The Annals of Statistics}, 19(3):1614--1625, 1991. \href{https://doi.org/10.1214/aos/1176348265}{doi:10.1214/aos/1176348265}.
\bibitem{giannelli2024} Carlotta Giannelli, Sofia Imperatore, Lisa Maria Kreusser, Estefan\'ia Loayza-Romero, Fatemeh Mohammadi, and Nelly Villamizar. A general formulation of reweighted least squares fitting. \emph{Mathematics and Computers in Simulation}, 225:52--65, 2024. \href{https://doi.org/10.1016/j.matcom.2024.04.029}{doi:10.1016/j.matcom.2024.04.029}.
\bibitem{saad2003} Yousef Saad. \emph{Iterative Methods for Sparse Linear Systems}. Second edition. Society for Industrial and Applied Mathematics, 2003. Author's electronic version: \url{https://www-users.cse.umn.edu/~saad/IterMethBook_2ndEd.pdf}.
\bibitem{alahmadi1993} Abdulhadi M. Alahmadi. \emph{Algorithms for the Construction of Constrained and Unconstrained Optimal Designs}. PhD thesis, University of Glasgow, July 1993. \url{https://theses.gla.ac.uk/76412/1/13834192.pdf}.
\bibitem{bertsekas2000} Dimitri P. Bertsekas and John N. Tsitsiklis. Gradient convergence in gradient methods with errors. \emph{SIAM Journal on Optimization}, 10(3):627--642, 2000. \url{https://www.mit.edu/~jnt/Papers/J079-00-gradient.pdf}.
\bibitem{razaviyayn2013} Meisam Razaviyayn, Mingyi Hong, and Zhi-Quan Luo. A Unified Convergence Analysis of Block Successive Minimization Methods for Nonsmooth Optimization. \emph{SIAM Journal on Optimization}, 23(2):1126--1153, 2013. \href{https://doi.org/10.1137/120891009}{doi:10.1137/120891009}. The detailed SUM assumptions and Theorem~1 cited here refer to \href{https://arxiv.org/abs/1209.2385v1}{arXiv:1209.2385v1}, 2012.
\bibitem{mairal2015} Julien Mairal. Incremental Majorization-Minimization Optimization with Application to Large-Scale Machine Learning. \emph{SIAM Journal on Optimization}, 25(2):829--855, 2015. \href{https://doi.org/10.1137/140957639}{doi:10.1137/140957639}. \url{https://lear.inrialpes.fr/people/mairal/resources/pdf/95763.pdf}.
\end{thebibliography}

\end{document}
